\subsection{FUNDAMENTAL SYSTEMS OF INTEGRALS OF A LINEAR SYSTEM OF SCALAR DIFFERENTIAL EQUATIONS}

In this subsection and the next we shall consider the case where E is a vector space of finite dimension n over the field C of complex numbers (so of dimension 2n over R), and where, for every \(t \in J\) , \(A(t)\) is an endomorphism of E with respect to the vector space structure over C. One can then identify \(A(t)\) with its matrix \((a_{ij}(t))\) with respect to a basis of E (over the field C), the \(a_{ij}\) this time being \(n^{2}\) complex functions defined and regulated on J; letting \(x_{j}\) ( \(1 \leqslant j \leqslant n\) ) denote the (complex) components of a vector \(x \in E\) with respect to the chosen basis, the linear equation

\[
{\frac {d \mathbf {x}}{d t}} = A (t). \mathbf {x} + \mathbf {b} (t)\tag{2}
\]

is again equivalent to the system

\[
\frac {d x _ {i}}{d t} = \sum_ {j = 1} ^ {n} a _ {i j} (t) x _ {j} + b _ {i} (t) \quad (1 \leqslant i \leqslant n).\tag{3}
\]

Theorems 1 (IV, p.~179) and 2 (IV, p.~179) and prop. 2 (IV, p.~181) then show that for every \(\mathbf{x}_0 = (x_{k0})_{1\leqslant k\leqslant n}\) in E there exists one and only one integral \(\mathbf{u} = (u_k)_{1\leqslant k\leqslant n}\) of the equation

\[
{\frac {d \mathbf {x}}{d t}} = A (t). \mathbf {x}\tag{4}
\]

defined on E and equal to \(x_{0}\) at the point \(t_{0}\) ; this integral can be written

\[
\mathbf {u} (t, t _ {0}, \mathbf {x} _ {0}) = C (t, t _ {0}). \mathbf {x} _ {0},
\]

\(C(t, t_{0})\) being an invertible square matrix \((c_{ij}(t, t_{0}))\) of order n whose entries are continuous complex functions on J × J and such that \(t \mapsto c_{ij}(t, t_{0})\) is a primitive of a regulated function on J.

In the particular case where \(n = 1\) the system (3) reduces to a single scalar equation

\[
{\frac {d x}{d t}} = a (t) x + b (t)\tag{11}
\]

\((a(t) \text{ and } b(t) \text{ being complex regulated functions on J}); \text{ one verifies immediately that the (one element) matrix } C(t, t_0) \text{ is equal to } \exp\left(\int_{t_0}^{t} a(s) ds\right); \text{ the integral of (11) equal to } x_0 \text{ at the point } t_0 \text{ is thus given explicitly by the formula}\)

\[
u (t) = x _ {0} \exp \left(\int_ {t _ {0}} ^ {t} a (s) d s\right) + \int_ {t _ {0}} ^ {t} b (s) \exp \left(\int_ {t _ {0}} ^ {t} a (\tau) d \tau\right) d s.\tag{12}
\]

In the space \(\mathcal{C}(\mathrm{J};\mathrm{E})\) of continuous maps from J into E, endowed with the topology of compact convergence, the set I of integrals of equation (4) is a vector subspace (over C) isomorphic to E, therefore to \(C^{n}\) (IV, p.~180, cor. 1, and IV, p.~181, prop. 2).

A basis \((\mathbf{u}_{j})_{1\leqslant j\leqslant n}\) of this space (over the field C) is called a fundamental system of integrals of (4).

PROPOSITION 4. For the n integrals \(\mathbf{u}_{j}\ (1 \leqslant j \leqslant n)\) of equation (4) to form a fundamental system it is necessary and sufficient that their values \(\mathbf{u}_{j}(t_{0})\) at a point \(t_{0} \in J\) be linearly independent vectors in E.

Indeed, the map which to every \(\mathbf{x}_0 \in \mathrm{E}\) associates the integral \(t \mapsto C(t, t_0). \mathbf{x}_0\) is an isomorphism of \(\mathrm{E}\) onto \(\mathcal{I}\) (IV, p.~180, cor. 1 and IV, p.~181, prop. 2).

If \(\left(\mathbf{e}_j\right)_{1\leqslant j\leqslant n}\) is any basis of E over C, the \(n\) integrals

\[
\mathbf {u} _ {j} (t) = C \left(t, t _ {0}\right). \mathbf {e} _ {j} \quad (1 \leqslant j \leqslant n)
\]

thus form a fundamental system; if one identifies \(C(t, t_0)\) with its matrix with respect to the basis \((\mathbf{e}_j)\) the integrals \(\mathbf{u}_j\) are precisely the columns of the matrix \(C(t, t_0)\) . The integral of (4) that takes the value \(\mathbf{x}_0 = \sum_{j=1}^{n} \lambda_j \mathbf{e}_j\) at the point \(t_0\) is then \(C(t, t_0). \mathbf{x}_0 = \sum_{k=1}^{n} \lambda_k \mathbf{u}_k(t)\) .

Given any \(n\) integrals \(\mathbf{u}_j\) ( \(1 \leqslant j \leqslant n\) ) of (4) one terms the determinant of these \(n\) integrals at a point \(t \in \mathbf{J}\) with respect to a basis \((\mathbf{e}_j)_{1 \leqslant j \leqslant n}\) of \(E\) , the determinant

\[
\Delta (t) = \left(\mathbf {u} _ {1} (t), \mathbf {u} _ {2} (t), \dots , \mathbf {u} _ {n} (t)\right)\tag{13}
\]

of the \(n\) vectors \(\mathbf{u}_j(t)\) with respect to the basis \((\mathbf{e}_j)\) (Alg., III, p.~522). One has (Alg., III, p.~523, prop. 2)

\[
\Delta (t) = \Delta (t _ {0}) \det \left(C (t, t _ {0})\right).\tag{14}
\]

By prop. 4 of IV, p.~184, for \(\left(\mathbf{u}_{j}\right)_{1\leqslant j\leqslant n}\) to be a fundamental system of integrals of (4) it is necessary and sufficient that the determinant \(\Delta(t)\) of the \(u_{j}\) be \(\neq0\) at some one point \(t_{0}\) of J; the formula (14) then shows that \(\Delta(t)\neq0\) at every point of J, in other words, that the vectors \(\mathbf{u}_{j}(t)\) ( \(1\leqslant j\leqslant n\) ) are always linearly independent.

PROPOSITION 5. The determinant of the matrix \(C(t, t_{0})\) is given by the formula

\[
\det \left(C (t, t _ {0})\right) = \exp \left(\int_ {t _ {0}} ^ {t} \operatorname{Tr} (A (s)) d s\right).\tag{15}
\]

Indeed, if one puts \(\delta(t) = \det\left(C(t, t_0)\right)\) one has, by the formula for the derivative of a determinant (I, p.~8, formula (3))

\[
\frac {d \delta}{d t} = \operatorname{Tr} \left(\frac {d C (t , t _ {0})}{d t} \left(C (t, t _ {0})\right) ^ {- 1}\right) \delta (t)
\]

that is, by the differential equation (5) of IV, p.~179 satisfied by \(C(t, t_0)\) ,

\[
\frac {d \delta}{d t} = \mathrm{Tr} \big (A (t) \big) \delta (t).
\]

Since \(\delta(t_0) = 1\) the formula (15) follows from the expression (12) (IV, p.~183) for the integral of a scalar linear equation.

Specifying n linearly independent integrals of (4) determines all the integrals of this equation, as we have just seen. We shall now show that for \(1 \leqslant p \leqslant n\) the knowledge of p linearly independent integrals \(\mathbf{u}_{j}\) ( \(1 \leqslant j \leqslant p\) ) of equation (4) reduces the integration of this equation to that of a homogeneous linear system of n - p scalar equations. Suppose that on an interval \(K \subset J\) there are n - p maps

\[
\mathbf {u} _ {p + k} \quad (1 \leqslant k \leqslant n - p)
\]

of K into E, which are primitives of regulated functions on K, and such that, for every \(t \in K\) , the n vectors \(\mathbf{u}_{j}(t) (1 \leqslant j \leqslant n)\) form a basis of E.

For every point \(t_1 \in \mathbf{J}\) there always exists an interval \(\mathbf{K}\) , a neighbourhood of \(t_1\) in \(\mathbf{J}\) , on which there are defined \(n - p\) functions \(\mathbf{u}_{p+k}\) ( \(1 \leqslant k \leqslant n - p\) ) having the preceding properties. For, let \((\mathbf{e}_i)_{1 \leqslant i \leqslant n}\) be a basis of \(\mathbf{E}\) ; there exist \(n - p\) vectors of this basis which form with the \(\mathbf{u}_j(t_1)\) ( \(1 \leqslant j \leqslant p\) ) a basis of \(\mathbf{E}\) (Alg., II, p.~292, th. 2); suppose for example that they are \(\mathbf{e}_{p+1}, \ldots, \mathbf{e}_n\) ; since the determinant \(\det(\mathbf{u}_1(t), \ldots, \mathbf{u}_p(t), \mathbf{e}_{p+1}, \ldots, \mathbf{e}_n)\) (with respect to the basis \((\mathbf{e}_i)\) ) is a continuous function of \(t\) and does not vanish for \(t = t_1\) there exists a neighbourhood \(\mathbf{K}\) of \(t_1\) on which it does not vanish; one can then take \(\mathbf{u}_{p+k}(t) = \mathbf{e}_{p+k}\) ( \(1 \leqslant k \leqslant n - p\) ) for \(t \in \mathbf{K}\) .

There exists an invertible matrix \(B(t)\) of order n, whose elements are primitives of regulated functions on K, such that \(B(t).e_{j} = \mathbf{u}_{j}(t)\) for \(1 \leqslant j \leqslant n\) . Let us put \(\mathbf{x} = B(t).y\) ; then y satisfies the equation \(\frac{dB}{dt}.y + B(t).\frac{dy}{dt} = A(t)B(t).y\) , which can also be written

\[
\frac {d \mathbf {y}}{d t} = (B (t)) ^ {- 1} \left(A (t) B (t) - \frac {d B}{d t}\right). \mathbf {y} = H (t). \mathbf {y}
\]

where \(H(t) = \left(h_{jk}(t)\right)\) is a matrix with regulated entries on K. By the definition of \(B(t)\) this linear equation admits the p constant vectors \(\mathbf{e}_{j}\) ( \(1 \leqslant j \leqslant p\) ) as integrals; one concludes immediately that necessarily \(h_{jk}(t) = 0\) for \(1 \leqslant k \leqslant p\) ; thus the components \(y_{k}\) of y (with respect to the basis \((\mathbf{e}_{i})\) ) with index \(k \geqslant p + 1\) satisfy a homogeneous linear system of n - p equations; once the solutions of this system are determined, the \(dy_{j}/dt\) for indices \(j \leqslant p\) are linear functions of the \(y_{k}\) with \(k \geqslant p + 1\) , so are known, and the primitives of these functions will give the \(y_{j}\) for indices \(j \leqslant p\) .

In particular, when one knows \(n - 1\) linearly independent integrals of equation (4) of IV, p.~183, integrating this equation reduces to that of a single homogeneous scalar equation, and then the evaluation of \(n\) primitives.

Remarks. 1) All the above applies also to the case where E is of dimension n over the field R and \(A(t)\) is an endomorphism of E for every \(t \in J\) : one has only to replace C by R throughout.

\begin{enumerate}
\def\labelenumi{\arabic{enumi})}
\setcounter{enumi}{1}
\tightlist
\item
  Let \(A(t) = (a_{ij}(t))\) be a square matrix of order n whose entries are real (resp. complex) regulated functions of t on J, and let \(C(t, t_0) = (c_{ij}(t, t_0))\) be the resolvent matrix of the corresponding linear system (3) (IV, p.~22). Let F be an arbitrary complete normed space over R (resp. C) and let us consider the system of linear differential equations
\end{enumerate}

\[
\frac {d \mathbf {y} _ {i}}{d t} = \sum_ {j = 1} ^ {n} a _ {i j} (t) \mathbf {y} _ {j}
\]

where the unknown functions \(\mathbf{y}_j\) take their values in F. It is immediate that the solution \((\mathbf{u}_j)_{1\leqslant j\leqslant n}\) of this system such that \(\mathbf{u}_j(t_0) = \mathbf{d}_j\) for \(1\leqslant j\leqslant n\) ( \(\mathbf{d}_j\) being arbitrary in F) is given by the formulae

\[
\mathbf {u} _ {i} (t) = \sum_ {j = 1} ^ {n} c _ {i j} (t, t _ {0}) \mathbf {d} _ {j} \quad (1 \leqslant i \leqslant n).
\]

We consider in particular the case where \(A(t)\) is an endomorphism of a vector space E of finite dimension n over C, such that there exists a basis of E with respect to which the matrix of \(A(t)\) has real elements for every \(t \in J\) . Then the above shows (by th. 1 of IV, p.~179) that the resolvent matrix \(C(t, t_{0})\) with respect to the same basis also has real elements: it suffices to consider the vector space \(E_{0}\) over R generated by the basis of E under consideration, and to remark that the restriction of \(A(t)\) to \(E_{0}\) is an endomorphism of this vector space.

\subsection{ADJOINT EQUATION}

Assuming always that the space E is of finite dimension n over C, let \(E^{*}\) be its dual (A, II, p.~40), which is a space of dimension n over C (Alg., II, p.~299, th. 4); the canonical bilinear form \(\langle x, x^{*}\rangle\) defined on \(E \times E^{*}\) (Alg., II, p.~234) is continuous on this product (being a polynomial in the components of \(x \in E\) and \(x^{*} \in E^{*}\) ).

Given a homogeneous linear equation (4) (IV, p.~183), where \(t \mapsto A(t)\) is a regulated map of \(\mathbf{J}\) into \(\mathcal{L}(\mathrm{E})\) , let us see if there exists a map \(t \mapsto \mathbf{v}(t)\) of \(\mathbf{J}\) into \(\mathrm{E}^*\) , a primitive of a regulated function on \(\mathbf{J}\) , and such that the scalar function \(t \mapsto \langle \mathbf{u}(t), \mathbf{v}(t) \rangle\) is constant on \(\mathbf{J}\) when \(\mathbf{u}\) is an arbitrary solution of (4); it comes to the same to write that the derivative of this function should be zero at every point where \(\mathbf{u}\) and \(\mathbf{v}\) are differentiable, that is, one must have

\[
\left\langle \frac {d {\bf u}}{d t}, {\bf v} (t) \right\rangle + \left\langle {\bf u} (t), \frac {d {\bf v}}{d t} \right\rangle = 0
\]

at such points.

Now, by (4), \(\left\langle\frac{d\mathbf{u}}{dt},\mathbf{v}(t)\right\rangle=\langle A(t).\mathbf{u}(t),\mathbf{v}(t)\rangle=-\langle\mathbf{u}(t),B(t).\mathbf{v}(t)\rangle\) where \(-B(t)\) is the transpose of \(A(t)\) (Alg., II, p.~234). The relation that v must satisfy can thus be written

\[
\left\langle \mathbf {u} (t), \frac {d \mathbf {v}}{d t} - B (t). \mathbf {v} (t) \right\rangle = 0
\]

at all points where \(A(t)\) is continuous and \(\mathbf{v}(t)\) is differentiable. Now for such a point \(t\) and an arbitrary point \(\mathbf{x}_0 \in E\) there exists, by th. 1 of IV, p.~179, a solution \(\mathbf{u}\) of (4) such that \(\mathbf{u}(t) = \mathbf{x}_0\) ; one thus must have \(\left\langle \mathbf{x}_0, \frac{d\mathbf{v}}{dt} - B(t).\mathbf{v}(t) \right\rangle = 0\) for all \(\mathbf{x}_0 \in E\) , which entails \(\frac{d\mathbf{v}}{dt} - B(t).\mathbf{v}(t) = 0\) . Consequently:

PROPOSITION 6. The map \(t \mapsto \mathbf{v}(t)\) of \(\mathbf{J}\) into \(\mathrm{E}^*\) , a primitive of a regulated function on \(\mathbf{J}\) , is such that \(\langle \mathbf{u}(t), \mathbf{v}(t) \rangle\) is constant on \(\mathbf{J}\) for every solution \(\mathbf{u}\) of the equation (4) of IV, p.~183 if and only if \(\mathbf{v}\) is a solution of the homogeneous linear equation

\[
{\frac {d \mathbf {x}}{d t}} = B (t). \mathbf {x}\tag{16}
\]

where \(-B(t)\) is the transpose of \(A(t)\) .

Equation (16) is called the adjoint of (4); clearly (4) is the adjoint of (16). The elements of the matrix \(B(t)\) being regulated functions of t on J, the results obtained above on linear equations apply to (16). In particular, the integral of (16) taking the value \(x_{0}^{*}\) at the point \(t_{0}\) can be written as \(H(t, t_{0})\) . \(x_{0}^{*}\) , where \(H(t, t_{0})\) is a bijective linear map of \(E^{*}\) onto itself, identical to the integral of the equation

\[
\frac {d V}{d t} = B (t) V\tag{17}
\]

which takes the value \(I\) at the point \(t_0\) . As a result one has (with the notation of IV, p.~179)

\[
\left\langle C (t, t _ {0}). \mathbf {x} _ {0}, H (t, t _ {0}). \mathbf {x} _ {0} ^ {*} \right\rangle = \left\langle \mathbf {x} _ {0}, \mathbf {x} _ {0} ^ {*} \right\rangle
\]

for any \(x_{0} \in E\) and \(x_{0}^{*} \in E^{*}\) , which shows that

\[
H (t, t _ {0}) = \check {C} (t, t _ {0})\tag{18}
\]

(the contragradient of \(C(t, t_{0})\) ). In particular, if one knows a fundamental system of integrals of the adjoint equation (16) the matrix \(H(t, t_{0})\) is determined, so is \(C(t, t_{0})\) also, and as a result, so are all the integrals of equation (4).

Remark. Let E and F be two complete normed spaces over R (or over C), and \((\mathbf{x}, \mathbf{y}) \mapsto \langle \mathbf{x}, \mathbf{y} \rangle\) a continuous bilinear form on \(E \times F\) , such that the relation `` \(\langle x, y \rangle = 0\) for all \(y \in F\) '' (resp. \(\langle x, y \rangle = 0\) for all \(x \in E\) '') implies x = 0 (resp. y = 0). Suppose further that for every \(t \in J\) there is a continuous linear map \(B(t)\) of F into itself, such that \(\langle A(t).x, y \rangle + \langle x, B(t).y \rangle = 0\) for every \((\mathbf{x}, \mathbf{y}) \in E \times F\) . In these circumstances one sees as before that for a map \(t \mapsto v(t)\) of J into F, a primitive of a regulated function, to be such that \(\langle u(t), v(t) \rangle\) is constant for every integral u of (4), it is necessary and sufficient that v should be an integral of (16), which again we call the adjoint of (4).

\subsection{LINEAR DIFFERENTIAL EQUATIONS WITH CONSTANT COEFFICIENTS}

Suppose again that E is an arbitrary complete normed space over R; let A be a continuous endomorphism of A, independent of t, and consider the homogeneous linear equation

\[
{\frac {d \mathbf {x}}{d t}} = A. \mathbf {x}.\tag{19}
\]

When \(\mathrm{E}\) is of finite dimension the equation (19) is equivalent to a homogeneous system (3) (IV, p.~183) of scalar differential equations, where the coefficients \(a_{ij}\) are constant.

By th. 1 (IV, p.~179), every integral of (19) is defined on all of \(\mathbf{R}\) ; by th. 2 (IV, p.~179) the integral of (19) taking the value \(\mathbf{x}_0\) at a point \(t_0 \in \mathbf{R}\) can be written as \(C(t, t_0)\mathbf{x}_0\) , where \(C(t, t_0)\) is a bijective bicontinuous linear map of E onto itself that satisfies the equation

\[
\frac {d U}{d t} = A U\tag{20}
\]

and is such that \(C(t_0, t_0) = I\) . Moreover, we have the identity

\[
C (t + \tau , t _ {0} + \tau) = C (t, t _ {0})\tag{21}
\]

for any \(\tau \in \mathbf{R}\) : indeed, one has \(dC(s, t_0 + \tau) / ds = AC(s, t_0 + \tau)\) by (20), and since \(A\) is constant one also has

\[
\frac {d C (t + \tau , t _ {0} + \tau)}{d t} = A C (t + \tau , t _ {0} + \tau);
\]

moreover

\[
C (t _ {0} + \tau , t _ {0} + \tau) = I = C (t _ {0}, t _ {0}),
\]

whence we have the identity (21), since the integral of (20) equal to I at the point \(t_{0}\) is unique.

If one puts \(C_0(t) = C(t, 0)\) then \(C(t, t_0) = C_0(t - t_0)\) ; moreover, for every \(\lambda \in \mathbf{R}\) , \(C_0(\lambda t)\) is identical to the integral of the equation

\[
\frac {d U}{d t} = \lambda A U\tag{22}
\]

which takes the value \(I\) at the point 0. We make the following definition:

DEFINITION 1. Given a continuous endomorphism A of E we denote by \(e^{A}\) or exp A the automorphism of E equal to the value at the point t = 1 of the integral of the equation (20) which takes the value I at the point t = 0.

With this notation the remarks preceding def. 1 show that

\[
C (t, t _ {0}) = \exp \left(A (t - t _ {0})\right).\tag{23}
\]

The exponential notation just introduced is justified by the following properties, which are entirely analogous to those of the function \(\exp z\) , for z real or complex (cf.~III, p.~98 and 106):

PROPOSITION 7. \(1^{\circ}\) The map \(X \mapsto e^{X}\) is a continuous map of \(\mathcal{L}(\mathrm{E})\) into the group of automorphisms of \(\mathrm{E}\) (invertible elements of \(\mathcal{L}(\mathrm{E})\) ).

\(2^{\circ}\) The map \(t\mapsto e^{Xt}\) of \(\mathbf{R}\) into \(\mathcal{L}(\mathrm{E})\) is differentiable, and

\[
\frac {d}{d t} \big (e ^ {X t} \big) = X e ^ {X t} = e ^ {X t} X.\tag{24}
\]

\(3^{\circ}\) For any \(X\in \mathcal{L}(\mathrm{E})\) one has

\[
e ^ {X} = \sum_ {n = 0} ^ {\infty} \frac {X ^ {n}}{n !}\tag{25}
\]

the right-hand side being absolutely and uniformly convergent on every bounded subset of \(\mathcal{L}(\mathrm{E})\) ; in particular, \(e^{It} = e^{t}I\) for \(t \in R\) .

\(4^{\circ}\) If X and Y commute then Y and \(e^{Y}\) commute with \(e^{X}\) , and

\[
e ^ {X + Y} = e ^ {X} e ^ {Y}.\tag{26}
\]

The relation (24) follows from the expression (23) for \(C(t,0)\) and from the fact that this function is an integral of (20); by recursion on n one deduces from (24) that \(t \mapsto e^{Xt}\) is indefinitely differentiable on R and that

\[
\mathrm{D} ^ {n} \left(e ^ {X t}\right) = X ^ {n} e ^ {X t}.
\]

By Taylor's formula one thus can write

\[
e ^ {X} = I + \frac {X}{1 !} + \frac {X ^ {2}}{2 !} + \dots + \frac {X ^ {n}}{n !} + X ^ {n + 1} \int_ {0} ^ {1} \frac {(1 - t) ^ {n}}{n !} e ^ {X t} d t.\tag{27}
\]

On the other hand, cor. 3 of IV, p.~181 shows that \(\left\| e^{Xt}\right\| \leqslant \exp \left(\| X\| |t|\right)\) . Thus the remainder \(r_n(X) = X^{n + 1}\int_0^1\frac{(1 - t)^n}{n!} e^{Xt}dt\) in formula (27) satisfies the inequality

\[
\| r _ {n} (X) \| \leqslant \frac {\| X \| ^ {n + 1}}{(n + 1) !} e ^ {\| X \|}
\]

whence one deduces the formula (25), the series on the right-hand side being absolutely and uniformly convergent on every bounded subset of \(\mathcal{L}(\mathrm{E})\) . For every pair of elements X, T of \(\mathcal{L}(\mathrm{E})\) one thus has

\[
e ^ {X + T} - e ^ {X} = \sum_ {n = 1} ^ {\infty} \frac {1}{n !} \big ((X + T) ^ {n} - X ^ {n} \big).
\]

Now, one can write \((X + T)^n - X^n = \sum_{(V_i)} V_1 V_2 \ldots V_n\) , where the sum is taken over the \(2^n - 1\) sequences \((V_i)\) of elements of \(\mathcal{L}(\mathrm{E})\) such that \(V_i = X\) or \(V_i = T\) for \(1 \leqslant i \leqslant n\) , and at least one of the \(V_i\) is equal to \(T\) ; the inequality

\[
\left\| (X + T) ^ {n} - X ^ {n} \right\| \leqslant \left(\| X \| + \| T \|\right) ^ {n} - \| X \| ^ {n}
\]

follows immediately, whence

\[
\| \exp (X + T) - \exp X \| \leqslant \exp \left(\| X \| + \| T \|\right) - \exp \| X \|\tag{28}
\]

which establishes the continuity of the map \(X \mapsto \exp X\) .

Finally, if X and Y commute, then Y commutes with \(e^{Xt}\) (IV, p.~181, prop. 2), so

\[
\frac {d}{d t} \big (e ^ {X t} e ^ {Y t} \big) = X e ^ {X t} e ^ {Y t} + e ^ {X t} (Y e ^ {Y t}) = (X + Y) e ^ {X t} e ^ {Y t}.
\]

Since, on the other hand, \(e^{Xt}e^{Yt}\) is equal to I for t = 0, one has \(e^{Xt}e^{Yt} = e^{(X+Y)t}\) , whence formula (26). From this latter, one deduces in particular that for arbitrary real s and t one has

\[
e ^ {X (s + t)} = e ^ {X s} e ^ {X t}\tag{29}
\]

and also that

\[
e ^ {- X} = \left(e ^ {X}\right) ^ {- 1}.\tag{30}
\]

On the contrary one notes that (26) need not remain valid if X and Y are no longer assumed to commute: indeed it would imply that \(\exp X\) and \(\exp Y\) always commute, which is not the case, as is shown by simple examples (IV, p.~204, exerc. 3).

Now let us assume that E is a vector space of finite dimension over the field C, and A an endomorphism of E (for its vector space structure over C) which one can identify with its matrix with respect to a basis of E; then, for all \(t \in R\) , \(e^{At}\) is an automorphism of E for the same structure (IV, p.~181, prop. 2). Let \(r_k\) ( \(1 \leqslant k \leqslant q\) ) be the distinct roots (in C) of the characteristic polynomial \(\varphi(r) = \det(A - rI)\) of the endomorphism A (the ``characteristic roots'' of A); if \(n_k\) is the order of multiplicity of \(r_k\) then \(\sum_{k=1}^{q} n_k = n\) . One knows that (Alg., VII, 31, n° 3) to each root \(r_k\) there corresponds a subspace \(E_k\) of E, of dimension \(n_k\) , such that \(E_k\) is stable under A, and that E is the direct sum of the \(E_k\) : \(E_k\) can be defined as the subspace of vectors x such that

\[
(A - r _ {k} I) ^ {n _ {k}}. \mathbf {x} = 0.
\]

Let \(\mathbf{a}\) be any vector in E; one can write \(\mathbf{a} = \sum_{k=1}^{q} \mathbf{a}_k\) , where \(\mathbf{a}_k \in \mathrm{E}_k\) ; the integral of the equation (19) of IV, p.~188, taking the value \(\mathbf{a}\) at the point \(t = 0\) is thus given by

\[
\mathbf {u} (t) = e ^ {A t}. \mathbf {a} = \sum_ {k = 1} ^ {q} e ^ {A t}. \mathbf {a} _ {k} = \sum_ {k = 1} ^ {q} e ^ {r _ {k} t} e ^ {(A - r _ {k} I) t}. \mathbf {a} _ {k}.
\]

But since \(\mathbf{a}_k\in \mathrm{E}_k\) one has

\[
\begin{array}{c} e ^ {(A - r _ {k} I) t}. \mathbf {a} _ {k} = \mathbf {a} _ {k} + \frac {t}{1 !} (A - r _ {k} I). \mathbf {a} _ {k} + \frac {t ^ {2}}{2 !} (A - r _ {k} I) ^ {2}. \mathbf {a} _ {k} + \dots \\ + \frac {t ^ {n _ {k} - 1}}{(n _ {k} - 1) !} (A - r _ {k} I) ^ {n _ {k} - 1}. \mathbf {a} _ {k}. \end{array}\tag{31}
\]

Thus every integral of the equation (19) of IV, p.~188, can be written as

\[
\mathbf {u} (t) = \sum_ {k = 1} ^ {q} e ^ {r _ {k} t} \mathbf {p} _ {k} (t)\tag{32}
\]

where \(\mathbf{p}_{k}(t)\) is a polynomial in t, with coefficients in the vector space \(E_{k}\) , and of degree \(\leqslant n_{k}-1\) . In particular, if all the roots of the characteristic equation of A are simple, the spaces \(E_{k}\) ( \(1 \leqslant k \leqslant n\) ) are all of dimension 1 over the field C, so there exist n vectors \(c_{k}\) such that the n functions \(e^{r_{k}t}c_{k}\) ( \(1 \leqslant k \leqslant n\) ) form a fundamental system of integrals of the equation (19) of IV, p.~188.

The characteristic roots of the endomorphism \(A\) are also called characteristic roots of the linear equation (19) of IV, p.~188. One may observe that one obtains the characteristic equation of \(A\) by writing that the function \(\mathbf{c}e^{rt}\) is an integral of (19) for a vector \(\mathbf{c} \neq 0\) .

When one has determined the roots \(r_{k}\) ( \(1 \leqslant k \leqslant q\) ) explicitly, and thus the order of multiplicity \(n_{k}\) of \(r_{k}\) , in practice one obtains the integrals of (19) by writing that this equation is satisfied by the expression (32) of IV, p.~191, where \(p_{k}\) is an arbitrary polynomial of degree \(\leqslant n_{k}-1\) , with coefficients in E; on identifying the coefficients of \(e^{r_{k}t}\) (for \(1 \leqslant k \leqslant q\) ) in the two sides of the equation so obtained one obtains linear equations for the coefficients of the polynomials \(p_{k}\) : one establishes easily that these equations determine the terms of degree \textgreater0 of \(p_{k}\) as functions of the constant term, and that the latter is a solution of the equation \((A - r_{k}I)^{n_{k}}.x = 0\) , which defines the subspace \(E_{k}\) (method of ``undetermined coefficients'').

Remark. When there exists a basis of E such that the matrix of A with respect to this basis has real elements (cf.~IV, p.~186, Remark 2), the characteristic equation of A has real coefficients. For every \(\mathbf{x} = (\xi_k)_{1 \leqslant k \leqslant n}\) of E, expressed in the basis under consideration, let \(\overline{x} = (\overline{\xi}_k)_{1 \leqslant k \leqslant n}\) ; the map \(x \mapsto \overline{x}\) is an antilinear involution of E. One knows (Alg., VII) that, if \(r_k\) is a non-real root of the characteristic equation, and \(E_k\) is the corresponding stable subspace, then \(\overline{r}_k\) is a characteristic root having the same multiplicity \(n_k\) as \(r_k\) , and the image \(E'_k\) of \(E_k\) under the map \(x \mapsto \overline{x}\) is the stable subspace corresponding to \(\overline{r}_k\) . One deduces from this that if \(u_j\) ( \(1 \leqslant j \leqslant n_k\) ) are \(n_k\) linearly independent integrals with values in \(E_k\) , then the \(2n_k\) integrals \(u_j + \overline{u}_j\) , \(i(\mathbf{u}_j - \overline{\mathbf{u}}_j)\) are linearly independent, and have, with respect to the chosen basis of E, components which are real functions of E. If \(r_k\) is a real characteristic root Remark 2 of IV, p.~186, shows that (with the same notation) there are \(n_k\) linearly independent integrals \(v_j\) ( \(1 \leqslant j \leqslant n_k\) ) with values in \(E_k\) whose components are real. That way one obtains a fundamental system of integrals of (19) whose components are all real.

\subsection{LINEAR EQUATIONS OF ORDER n}

One calls a linear differential equation of order n an equation of the form

\[
\mathrm{D} ^ {n} x - a _ {1} (t) \mathrm{D} ^ {n - 1} x - \dots - a _ {n - 1} (t) \mathrm{D} x - a _ {n} (t) x = b (t)\tag{33}
\]

where the \(a_{k}\) ( \(1 \leqslant k \leqslant n\) ) and b are real (complex) functions of the real variable t, defined on an interval J of R. The general procedure of IV, p.~164 shows that this equation is equivalent to the linear system of n first order equations

\[
\left\{ \begin{array}{l l} \frac {d x _ {k}}{d t} = x _ {k + 1} & (1 \leqslant k \leqslant n - 1) \\ \frac {d x _ {n}}{d t} = a _ {1} (t) x _ {n} + a _ {2} (t) x _ {n - 1} + \dots + a _ {n} (t) x _ {1} + b (t) \end{array} \right.\tag{34}
\]

that is to say, to the linear equation

\[
\frac {d \mathbf {x}}{d t} = A (t). \mathbf {x} + \mathbf {b} (t)\tag{35}
\]

where we have put \(\mathbf{x}=(x_{1},x_{2},\ldots,x_{n})\in\mathbf{C}^{n},\mathbf{b}(t)=(0,0,\ldots,0,b(t))\) , and where the matrix \(A(t)\) is defined by

\[
A (t) = \left( \begin{array}{c c c c c} 0 & 1 & 0 & \ldots & 0 \\ 0 & 0 & 1 & \ldots & 0 \\ \vdots & \vdots & \vdots & \ldots & \vdots \\ 0 & 0 & 0 & \ldots & 1 \\ a _ {n} (t) & a _ {n - 1} (t) & a _ {n - 2} (t) & \ldots & a _ {1} (t) \end{array} \right).
\]

The study of the linear equation of order n thus consists of applying the general results above to the particular linear equation (35). For every interval J where the functions \(a_{j}\) ( \(1 \leqslant j \leqslant n\) ) and b are regulated there exists one and only one function u, defined on J, having a continuous derivative of order n - 1, and a regulated derivative of order n on this interval (except at the points of a countable set), satisfying (33) on the complement of a countable subset of J, and such that

\[
u (t _ {0}) = x _ {0}, \qquad \mathrm{D} u (t _ {0}) = x _ {0} ^ {\prime}, \dots , D ^ {n - 1} u (t _ {0}) = x _ {0} ^ {(n - 1)}\tag{36}
\]

where \(t_0\) is an arbitrary point of \(J\) , and \(x_0, x_0', \ldots, x_0^{(n-1)}\) are \(n\) arbitrary complex numbers.

For the \(p\) integrals \(u_{j}\) ( \(1 \leqslant j \leqslant p\) ) of the homogeneous equation

\[
\mathrm{D} ^ {n} x - a _ {1} (t) \mathrm{D} ^ {n - 1} x - \dots - a _ {n - 1} (t) \mathrm{D} x - a _ {n} (t) x = 0\tag{37}
\]

associated with (33) to be linearly independent (in the space \(\mathcal{C}(\mathrm{J};\mathbf{C})\) of continuous maps from J into C, considered as a vector space over C), it is necessary and sufficient that the corresponding p integrals \(\mathbf{u}_{j}=(u_{j},\mathrm{D}u_{j},\ldots,\mathrm{D}^{n-1}u_{j})\) of the homogeneous equation \(d\mathbf{x}/dt=A(t)\mathbf{x}\) should be linearly independent (in the space \(\mathcal{C}(\mathrm{J};\mathbf{C}^{n})\) of continuous maps from J into \(C^{n}\) ). It is clear that this condition is necessary.

Conversely, if there are \(n\) complex constants \(\lambda_j\) , not all zero, such that \(\sum_{j=1}^{n} \lambda_j u_j(t) = 0\) identically on \(J\) , one deduces that \(\sum_{j=1}^{n} \lambda_j D^k u_j(t) = 0\) on \(J\) for every integer \(k\) such that \(1 \leqslant k \leqslant n - 1\) , which means that \(\sum_{j=1}^{n} \lambda_j \mathbf{u}_j(t) = 0\) on \(J\) .

Consequently (IV, p.~180, cor. 1)

PROPOSITION 8. The set of integrals of the homogeneous linear equation (37), defined on J, is a vector space of dimension n over the field C.

Given any n integrals \(u_{j}\) ( \(1 \leqslant j \leqslant n\) ) of equation (37) one calls the Wronskian of this system of integrals the determinant (with respect to the canonical basis of \(C^{n}\) ) of the corresponding n integrals \(u_{j}\) of the equation \(d\mathbf{x}/dt = A(t)\mathbf{x}\) , that is to say, the function

\[
\mathrm{W} (t) = \left| \begin{array}{c c c c} u _ {1} (t) & u _ {2} (t) & \ldots & u _ {n} (t) \\ \mathrm{D} u _ {1} (t) & \mathrm{D} u _ {2} (t) & \ldots & \mathrm{D} u _ {n} (t) \\ \vdots & \vdots & \ldots & \vdots \\ \mathrm{D} ^ {n - 1} u _ {1} (t) & \mathrm{D} ^ {n - 1} u _ {2} (t) & \ldots & \mathrm{D} ^ {n - 1} u _ {n} (t) \end{array} \right|.
\]

For the n integrals \(u_{j}\) to be linearly independent it is necessary and sufficient that \(\mathrm{W}(t) \neq 0\) on J; moreover, it suffices for this that \(\mathrm{W}(t_{0}) \neq 0\) at only one \(t_{0}\) of J (IV, p.~184, prop. 4); further, one has (IV, p.~184, prop. 5)

\[
\mathrm{W} (t) = \mathrm{W} (t _ {0}) \exp \left(\int_ {t _ {0}} ^ {t} a _ {1} (s) d s\right).\tag{38}
\]

We identify the resolvent \(C(t, t_{0})\) of equation (35) with its matrix with respect to the canonical basis of \(C^{n}\) ; the columns \(\mathbf{v}_{j}(t, t_{0})\) ( \(1 \leqslant j \leqslant n\) ) of this matrix are then n linearly independent integrals

\[
\mathbf {v} _ {j} (t, t _ {0}) = (v _ {j} (t, t _ {0}), \mathrm{D} v _ {j} (t, t _ {0}), \dots , \mathrm{D} ^ {n - 1} v _ {j} (t, t _ {0}))
\]

of the homogeneous equation \(d\mathbf{x}/dt = A(t)\) . x which correspond to n linearly independent integrals \(v_{j}(t, t_{0})\) of equation (37) such that

\[
\mathbf {D} ^ {k - 1} v _ {j} (t _ {0}, t _ {0}) = \delta_ {j k}
\]

(Kronecker delta) for \(1 \leqslant j \leqslant n\) , \(1 \leqslant k \leqslant n\) (on agreeing to put \(D^{0}v_{j} = v_{j}\) ). It results in particular that the method of variation of parameters (IV, p.~182) applied to equation (35) here gives as a particular integral of (33), equal to 0 together with its first n - 1 derivatives at the point \(t_{0}\) , the function

\[
w (t) = \int_ {t _ {0}} ^ {t} v _ {n} (t, s) b (s) d s.\tag{39}
\]

In the particular case of the equation \(D^{n}x = b(t)\) the formula (39) again gives the formula expressing the \(n^{th}\) primitive of the regulated function \(b(t)\) which vanishes with its first n - 1 derivatives at the point \(t_{0}\) , namely

\[
w (t) = \int_ {t _ {0}} ^ {t} b (s) \frac {(t - s) ^ {n - 1}}{(n - 1) !} d s
\]

(II, p.~62, formula (19)): the integral of \(D^{n}x = 0\) which vanishes together with its first n - 2 derivatives at the point \(t_{0}\) , and whose \(n - 1^{th}\) derivative is equal to 1 there, is actually the polynomial \((t - t_{0})^{n-1}/(n - 1)!\) .

\subsection{LINEAR EQUATIONS OF ORDER n WITH CONSTANT COEFFICIENTS}

If the coefficients \(a_{j}\) in equation (33) are constant, the corresponding matrix \(A\) is constant; the characteristic equation is obtained by writing that \(e^{rt}\) is a solution, which gives

\[
r ^ {n} - a _ {1} r ^ {n - 1} - \ldots - a _ {n - 1} r - a _ {n} = 0.\tag{40}
\]

Let \(r_j\) ( \(1 \leqslant j \leqslant q\) ) be the distinct roots of this equation, and \(n_j\) ( \(1 \leqslant j \leqslant q\) ) the multiplicity of the root \(r_j \left( \sum_{j=1}^{q} n_j = n \right)\) . By the results of IV, p.~188 to 194, to each root \(r_j\) there corresponds, for the homogeneous equation

\[
\mathrm{D} ^ {n} x - a _ {1} \mathrm{D} ^ {n - 1} x - \dots - a _ {n - 1} \mathrm{D} x - a _ {n} x = 0\tag{41}
\]

a system of \(n_j\) linearly independent integrals

\[
u _ {j k} (t) = e ^ {r _ {j} t} p _ {j k} (t),
\]

where \(p_{jk}\) is a polynomial (with complex coefficients) of degree \(\leqslant n_{j}-1\) ( \(1\leqslant k\leqslant n_{j}\) ); further, the n integrals \(u_{jk}\) ( \(1\leqslant j\leqslant q\) , \(1\leqslant k\leqslant n_{j}\) ) so obtained are linearly independent. It follows that the \(n_{j}\) polynomials \(p_{jk}\) ( \(1\leqslant k\leqslant n_{j}\) ) are linearly independent in the space of polynomials in t of degree \(\leqslant n_{j}-1\) , so form a basis (over C) of this space, since the latter is of dimension \(n_{j}\) . In other words:

PROPOSITION 9. Let \(r_j\) ( \(1 \leqslant j \leqslant q\) ) be the distinct roots of the characteristic equation (40), and let \(n_j\) be the multiplicity of the root \(r_j\) ( \(1 \leqslant j \leqslant q\) ). Then the \(n\) functions \(t^k e^{r_j t}\) ( \(1 \leqslant k \leqslant n_j\) , \(1 \leqslant j \leqslant q\) ) are linearly independent integrals of the homogeneous equation (41).

One can prove this result directly in the following way. It follows from equation (41) that the \(n^{th}\) derivative of every integral of this equation is differentiable on R, from which one deduces immediately, by induction on the integer m \textgreater{} n, that every integral of (41) admits a derivative of order m, that is, is indefinitely differentiable on R. Let D be the (non-topological) vector space over C of indefinitely differentiable complex-valued functions on R; the map \(x \mapsto Dx\) is an endomorphism of this space, and equation (41) can be written

\[
f (\mathrm{D}) x = 0\tag{42}
\]

\[
\text { where } f (\mathrm{D}) = \mathrm{D} ^ {n} - a _ {1} \mathrm{D} ^ {n - 1} - \dots - a _ {n - 1} \mathrm{D} - a _ {n} (\mathrm{Alg}., \mathrm{IV}, \mathrm{p}. 4).
\]

PROPOSITION 10. Let \(g\) and \(h\) be two relatively prime polynomials such that \(f = gh\) . The subspace of solutions of (42) is the direct sum of the subspaces of solutions of the two equations

\[
g (\mathrm{D}) x = 0, \quad h (\mathrm{D}) x = 0.
\]

Now, by the Bezout identity (Alg., VII. 2, th. 1) there exist two polynomials \(p(\mathrm{D})\) and \(q(\mathrm{D})\) such that \(p(\mathrm{D})g(\mathrm{D}) + q(\mathrm{D})h(\mathrm{D}) = 1\) . For every solution \(x\) of (42) one can write \(x = y + z\) , where \(y = p(\mathrm{D})g(\mathrm{D})x\) and \(z = q(\mathrm{D})h(\mathrm{D})x\) ; and then \(h(\mathrm{D})y = p(\mathrm{D})(f(\mathrm{D})x) = 0\) and \(g(\mathrm{D})z = q(\mathrm{D})(f(\mathrm{D})x) = 0\) . On the other hand, if both \(g(\mathrm{D})x = 0\) and \(h(\mathrm{D})x = 0\) , one deduces that

\[
x = p (\mathrm{D}) (g (\mathrm{D}) x) + q (\mathrm{D}) (h (\mathrm{D}) x) = 0,
\]

which completes the proof.

With the preceding notation one can write

\[
f (\mathrm{D}) = \prod_ {j = 1} ^ {q} (\mathrm{D} - r _ {j}) ^ {n _ {j}}
\]

and prop. 10, by induction on \(q\) , shows that the subspace of solutions of (42) is the direct sum of the subspaces of solutions of the \(q\) equations

\[
(\mathrm{D} - r _ {j}) ^ {n _ {j}} x = 0 \quad (1 \leqslant j \leqslant q).\tag{43}
\]

Now, for every complex number r one has

\[
\mathrm{D} (e ^ {r t} x) = e ^ {r t} (\mathrm{D} + r) x\tag{44}
\]

so (43) is equivalent to

\[
\mathrm{D} ^ {n _ {j}} \left(e ^ {- r _ {j} t} x\right) = 0
\]

and thus has as solutions the functions \(e^{r_{j}t} p_{j}(t)\) , where \(p_{j}\) runs through the set of polynomials of degree \(\leqslant n_{j}-1\) ; thus one recovers prop. 9 of IV, p.~194.

Assuming that the homogeneous equation (41) has been solved (that is, assuming that the characteristic roots have been found), one knows that the method of variation of parameters allows one to find the solutions of the inhomogeneous equation

\[
\mathrm{D} ^ {n} x - a _ {1} \mathrm{D} ^ {n - 1} x - \dots - a _ {n - 1} \mathrm{D} x - a _ {n} x = b (t)\tag{45}
\]

where \(b(t)\) is an arbitrary regulated function (IV, p.~182); we note that if \(b(t)\) is indefinitely differentiable on an interval \(J\) then all the integrals of (45) are indefinitely differentiable on J. In the particular case \(b(t) = e^{\alpha t} p(t)\) , where p is a polynomial (with complex coefficients) and \(\alpha\) is an arbitrary complex number, one obtains an integral of (45) more simply by the following method. Put \(x = e^{\alpha t} y\) ; the equation

\[
f (\mathrm{D}) x = e ^ {\alpha t} p (t)
\]

can, by (44), be written

\[
f (\alpha + \mathrm{D}) y = p (t)
\]

or, by Taylor's formula applied to the polynomial \(f(\mathbf{D})\) ,

\[
\frac {f ^ {(n)} (\alpha)}{n !} \mathrm{D} ^ {n} y + \frac {f ^ {(n - 1)} (\alpha)}{(n - 1) !} \mathrm{D} ^ {n - 1} y + \dots + \frac {f ^ {\prime} (\alpha)}{1 !} \mathrm{D} y + f (\alpha) y = p (t).\tag{46}
\]

Let \(m\) be the degree of the polynomial \(p(t) = \sum_{k=0}^{m} \lambda_k t^{m-k}\) ; if \(f(\alpha) \neq 0\) (that is, if \(\alpha\) is not a characteristic root), there exists one and only one polynomial \(u(t) = \sum_{k=0}^{m} c_k t^{m-k}\) of degree \(m\) which is a solution of (46), for the coefficients \(c_k\) are determined by the system of linear equations

\[
\begin{array}{r l} f (\alpha) c _ {k} + \binom {m - k + 1} {1} f ^ {\prime} (\alpha) c _ {k - 1} + \binom {m - k + 2} {2} f ^ {\prime \prime} (\alpha) c _ {k - 2} + \dots \\ & + \binom {m} {k} f ^ {(k)} (\alpha) c _ {0} = \lambda_ {k} \quad (0 \leqslant k \leqslant m) \end{array}
\]

which clearly admits one and only one solution. If, on the other hand, \(\alpha\) is a characteristic root, and h is its multiplicity, the preceding calculation shows that there is one and only one polynomial of degree m such that every solution of \(D^{h}y = v(t)\) is an integral; in other words, every polynomial solution of (46) is then of degree \(m + h\) (``resonance'').

\subsection{SYSTEMS OF LINEAR EQUATIONS WITH CONSTANT COEFFICIENTS}

With the notation of \(n^{\circ}\) 8, let us consider more generally a system of m differential equations of the form

\[
\sum_ {k = 1} ^ {n} p _ {j k} (\mathrm{D}) x _ {k} = b _ {j} (t) \quad (1 \leqslant j \leqslant m)\tag{47}
\]

where the unknowns \(x_{k}\) ( \(1 \leqslant k \leqslant n\) ) and the right-hand sides \(b_{j}\) ( \(1 \leqslant j \leqslant m\) ) are complex functions of the real variable t, and where the \(p_{jk}(D)\) are polynomials (of any degree) with constant (complex) coefficients in the differentiation operator D ( \(1 \leqslant j \leqslant m, 1 \leqslant k \leqslant n\) ).

Such systems are not of the same type as those considered in IV, p.~1164 (formula (5)), as the following example shows:

\[
\left\{ \begin{array}{c} \mathrm{D} x _ {1} = a (t) \\ \mathrm{D} ^ {2} x _ {1} + \mathrm{D} x _ {2} + x _ {2} = b (t). \end{array} \right.\tag{48}
\]

We shall confine ourselves to the case where the functions \(b_{j}(t)\) are indefinitely differentiable on the interval J, and we shall look only for solutions \((x_{k})_{1\leqslant k\leqslant n}\) which are indefinitely differentiable on J. On putting \(\mathbf{b}(t)=\left(b_{1}(t),\ldots,b_{m}(t)\right)\) , (a map from J into \(C^{m}\) ), and \(\mathbf{x}=\left(x_{1},x_{2},\ldots,x_{n}\right)\) , the system (47) can be written as

\[
P (\mathrm{D}) \mathbf {x} = \mathbf {b} (t)\tag{49}
\]

where \(P(\mathrm{D})\) is the matrix \((p_{jk}(\mathrm{D}))\) with m rows and n columns, whose coefficients belong to the ring C{[}D{]} of polynomials in D, with coefficients in C. Let \(f_{j}(\mathrm{D})\) ( \(1 \leqslant j \leqslant r \leqslant \operatorname{Min}(m, n)\) ) be the nonzero similarity invariants of the matrix \(P(\mathrm{D})\) ; one knows (Alg., VII, p.~32) that these are well-determined monic polynomials, such that \(f_{j}\) divides \(f_{j+1}\) for \(1 \leqslant j \leqslant r - 1\) (r being the rank of \(P(\mathrm{D})\) ); further, there are two square matrices \(U(\mathrm{D})\) and \(V(\mathrm{D})\) of order m and n respectively, invertible (in the rings of square matrices of order m and n respectively, with coefficients in the ring C{[}D{]} of polynomials in D with complex coefficients), and such that all the entries of the matrix \(Q(\mathrm{D}) = (q_{jk}(\mathrm{D})) = U(\mathrm{D}) P(\mathrm{D}) V(\mathrm{D})\) are zero, apart from the diagonal terms \(q_{jj}(\mathrm{D}) = f_{j}(\mathrm{D})\) for \(1 \leqslant j \leqslant r\) . Now we put \(\mathbf{y} = V^{-1}(\mathbf{D})\mathbf{x}\) ; the equation (49) is equivalent to the equation

\[
U (\mathrm{D}) \big (P (\mathrm{D}) \big (V (\mathrm{D}) \mathbf {y} \big) \big) = U (\mathrm{D}) \mathbf {b}
\]

that is, to

\[
Q (\mathrm{D}) \mathbf {y} = U (\mathrm{D}) \mathbf {b}\tag{50}
\]

since \(U(\mathbf{D})\) is invertible. Now, if \(\mathbf{y} = (y_1, y_2, \ldots, y_n)\) , and if

\[
U (\mathrm{D}) \mathbf {b} (t) = \left(c _ {1} (t), \dots , c _ {m} (t)\right),
\]

then equation (50) can be written

\[
f _ {j} (\mathrm{D}) y _ {j} = c _ {j} (t) \quad \text {   for   } 1 \leqslant j \leqslant r\tag{51}
\]

\[
0 = c _ {j} (t) \quad \text {   for   } r + 1 \leqslant j \leqslant m.\tag{52}
\]

The system thus does not admit any indefinitely differentiable solutions unless the conditions (52) are satisfied; the determination of the \(y_{j}\) for indices \(j \leqslant r\) then reduces to integrating the r linear differential equations with constant coefficients (51); the \(y_{j}\) for indices \textgreater{} r are arbitrary indefinitely differentiable functions. Once the solutions y of (50) have thus been determined, one deduces the solutions of (47) from the formula \(\mathbf{x} = V(\mathrm{D})\mathbf{y}\) .

Remarks. 1) Some of the polynomials \(f_{j}(\mathrm{D})\) may reduce to nonzero constants; the corresponding \(y_{j}\) are then completely determined.

\begin{enumerate}
\def\labelenumi{\arabic{enumi})}
\setcounter{enumi}{1}
\item
  When the \(b_{j}\) are all zero, that is, if the system (47) is homogeneous, the conditions (52) are always satisfied; if, further, \(r = n\) , one sees that the set of solutions of (47) is a vector space over \(\mathbf{C}\) , of dimension equal to the sum of the degrees of the \(f_{j}(\mathrm{D})\) , that is, to the degree of \(\det(P(\mathrm{D}))\) .
\item
  Given the polynomials \(p_{jk}(\mathbf{D})\) , a system (47) which admits solutions when the right-hand sides are indefinitely differentiable (or differentiable of a certain order) may not admit them when the right-hand sides are arbitrary regulated functions: this is shown by the example (48), which has no solution when \(a(t)\) is not a primitive. We shall not undertake here to seek for supplementary possibility conditions which enter when the right-hand sides are arbitrary regulated functions.
\end{enumerate}

\setcounter{exercisegroup}{0}
\refstepcounter{exercisegroup}
\section*{Exercises on \S\arabic{exercisegroup}}
\phantomsection
\addcontentsline{toc}{section}{Exercises on \protect\S\arabic{exercisegroup}}

\begin{enumerate}
\def\labelenumi{\arabic{enumi})}
\tightlist
\item
  \begin{enumerate}
  \def\labelenumii{\alph{enumii})}
  \tightlist
  \item
    With the notation of IV, p.~165, suppose that \(\mathbf{f}\) is uniformly continuous on \(\mathrm{J} \times \mathrm{S}\) . Prove prop. 3 of IV, p.~166 without using the axiom of choice. (Let \(\delta\) be such that the relations \(|t_2 - t_1| \leqslant \delta\) , \(\| \mathbf{x}_2 - \mathbf{x}_1 \| \leqslant \delta\) imply \(\| \mathbf{f}(t_2, \mathbf{x}_2) - \mathbf{f}(t_1, \mathbf{x}_1) \| \leqslant \varepsilon\) ; consider a subdivision of J into intervals of length \(\leqslant \inf(\delta, \delta / M)\) and define the approximate solution by induction.)
  \end{enumerate}
\end{enumerate}

\begin{enumerate}
\def\labelenumi{\alph{enumi})}
\setcounter{enumi}{1}
\tightlist
\item
  When E is finite dimensional and f is Lipschitz on \(I \times H\) prove prop. 3 of IV, p.~166 without using the axiom of choice. (Remark that for every \(\delta > 0\) there exists a finite number of points \(x_{i}\) of S such that every point of S is at a distance \(\leqslant \delta\) from one of the \(x_{i}\) ; then proceed as in a), considering the regulated functions \(t \mapsto \mathbf{f}(t, \mathbf{x}_{i})\) , of which there are only a finite number.)
\end{enumerate}

\begin{enumerate}
\def\labelenumi{\arabic{enumi})}
\setcounter{enumi}{1}
\item
  Given two numbers \(b > 0\) , \(M > 0\) and an arbitrary number \(\varepsilon > 0\) , give an example of a scalar differential equation \(x' = g(x)\) such that \(|g(x)| \leqslant M\) for \(|x| \leqslant b\) and which admits an integral \(x = u(t)\) continuous on the interval \(] - b/M - \varepsilon, b/M + \varepsilon[\) , but which does not have a finite limit at the point \(x = \frac{b}{M} + \varepsilon\) (define \(g\) so that the integral considered has a continuous derivative on \(] - b/M - \varepsilon, b/M + \varepsilon[\) , this derivative being equal to the constant \(M\) on \(] - b/M, b/M[\) ).
\item
  Let \(S \subset H\) be an open ball with centre \(\mathbf{x}_0\) and radius \(r\) , and \(\mathbf{f}\) a Lipschitz function for the constant \(k > 0\) on \(I \times S\) ; suppose further that \(t \mapsto \mathbf{f}(t, \mathbf{x}_0)\) is bounded on \(I\) , and denote by \(M_0\) the supremum of \(\| \mathbf{f}(t, \mathbf{x}_0) \|\) over \(t \in I\) . Show that for all \(t_0 \in I\) there is an integral \(\mathbf{u}\) of \(\mathbf{x}' = \mathbf{f}(t, \mathbf{x})\) , with values in \(S\) , equal to \(\mathbf{x}_0\) at the point \(t_0\) and defined on the intersection of \(I\) and the interval \(]t_0 - \lambda, t_0 + \lambda[\) , where
\end{enumerate}

\[
\lambda = \frac {1}{k} \log \left(1 + \frac {k z}{M _ {0}}\right)
\]

(note that one has \(\| \mathbf{u}(t) - \mathbf{x}_0\| \leqslant \mathbf{M}(t - t_0) + k\int_{t_0}^t\| \mathbf{u}(s) - \mathbf{x}_0\| ds\) for \(t > t_0\) ).

* 4) a) Let \(I = ]t_0 - a, t_0 + a[\) be an open interval in \(\mathbf{R}\) , let \(S\) be an open ball with centre \(\mathbf{x}_0\) and radius \(r\) in \(E\) , and \(\mathbf{f}\) a locally Lipschitz function on \(I \times S\) . Let \((s, z) \mapsto h(s, z)\) be a function \(\geqslant 0\) defined and continuous for \(0 \leqslant s \leqslant a\) and \(0 \leqslant z \leqslant r\) , such that for every \(s \in [0, a]\) the map \(z \mapsto h(s, z)\) is increasing; suppose that \(\| \mathbf{f}(t, \mathbf{x}) \| \leqslant h(|t - t_0|, \| \mathbf{x} - \mathbf{x}_0 \|)\) on \(I \times S\) . Let \(\varphi\) be a primitive of a regulated function on an interval \([0, \alpha]\) (with \(\alpha < a\) ), with values in \([0, r[\) , such that \(\varphi(0) = 0\) and \(\varphi'(s) > h(s, \varphi(s))\) on \([0, \alpha]\) except on a countable set. Show that the integral \(\mathbf{u}\) of \(\mathbf{x}' = \mathbf{f}(t, \mathbf{x})\) , equal to \(\mathbf{x}_0\) at the point \(t_0\) , is defined on \(]t_0 - \alpha, t_0 + \alpha[\) , and that \(\| \mathbf{u}(t) - \mathbf{x}_0 \| \leqslant \varphi(|t - t_0|)\) on this interval.

\begin{enumerate}
\def\labelenumi{\alph{enumi})}
\setcounter{enumi}{1}
\tightlist
\item
  Deduce from a) that if f is defined on I × E, and there is a function h defined, continuous, increasing and \textgreater{} 0 on {[}0, +∞{[}, such that \(\int_{0}^{+\infty} dz / h(z) = +\infty\) , and that \(\|\mathbf{f}(t, \mathbf{x})\| \leqslant h(\|x\|)\) , then every integral of \(\mathbf{x}' = \mathbf{f}(t, \mathbf{x})\) is defined on all of I.
\end{enumerate}

* 5) Let \(E\) be the complete normed space of sequences \(\mathbf{x} = (x_{n})\) of real numbers such that \(\lim_{n\to \infty}x_n = 0\) , normed by \(\| \mathbf{x}\| = \sup_n|x_n|\) . For every integer \(n\geqslant 0\) let \(\mathbf{e}_n\) be the sequence every term of which is zero, apart from the term with index \(n\) , which is equal to 1; the space \(E\) is the direct sum of the subspace \(V_{n}\) of dimension 2 generated by \(\mathbf{e}_n\) and \(\mathbf{e}_{n + 1}\) , and the closed subspace \(W_{n}\) generated by the \(\mathbf{e}_k\) with indices different from \(n\) and \(n + 1\) . Let \(\mathbf{f}_n\) be a continuous Lipschitz function on \(E\) , with values in \(E\) , constant on every equivalence class modulo \(W_{n}\) , equal to \(\mathbf{e}_{n + 1} - \mathbf{e}_n\) on the line joining \(\mathbf{e}_n\) to \(\mathbf{e}_{n + 1}\) , and equal to 0 on the intersection of \(V_{n}\) and the ball \(\| \mathbf{x}\| \leqslant \frac{1}{4}\) . On the other hand, for every integer \(n > 0\) let \(\varphi_{n}\) be a real function defined and continuous on the interval \(\left[\frac{1}{n + 1},\frac{1}{n}\right]\) , equal to 0 at the endpoints of this interval, and such that \(\int_{1 / n + 1}^{1 / n}\varphi_n(t)dt = 1\) . Let \(I\) be the interval {[}0, 1{]} in \(\mathbf{R}\) ; consider on \(I\times E\) the function \(\mathbf{f}\) with values in \(E\) , defined as follows: \(\mathbf{f}(0,\mathbf{x}) = 0\) for every \(\mathbf{x}\in E\) ; for \(\frac{1}{n + 1}\leqslant t\leqslant \frac{1}{n}\) , let \(\mathbf{f}(t,\mathbf{x}) = \varphi_n(t)\mathbf{f}_n(\mathbf{x})\) . Show that \(\mathbf{f}\) is continuous and locally Lipschitz on \(I\times E\) , but that there exists an integral \(\mathbf{u}\) of the equation \(\mathbf{x}' = \mathbf{f}(t,\mathbf{x})\) , defined and bounded on {[}0, 1{]}, equal to \(\mathbf{e}_n\) for \(t = 1 / n\) , and consequently having no limit as \(t\) tends to 0.

* 6) a) Let \(\mathrm{I} = [0, +\infty[\) ; if \(\mathbf{f}\) is Lipschitz on \(\mathrm{I} \times \mathrm{E}\) for a regulated function \(k(t) > 0\) such that the integral \(\int_0^{+\infty} k(t) dt\) converges, show that every integral of the equation \(\mathbf{x}' = \mathbf{f}(t, \mathbf{x})\) is defined on all of \(\mathrm{I}\) .

\begin{enumerate}
\def\labelenumi{\alph{enumi})}
\setcounter{enumi}{1}
\tightlist
\item
  If further the integral \(\int_{0}^{+\infty}\|\mathbf{f}(t,\mathbf{x}_{0})\|dt\) converges (for a certain point \(x_{0}\in E\) ), show that every integral of \(\mathbf{x}^{\prime}=\mathbf{f}(t,\mathbf{x})\) tends to a finite limit as t tends to \(+\infty\) (first show that every integral is bounded as t tends to \(+\infty\) ).
\end{enumerate}

\begin{enumerate}
\def\labelenumi{\arabic{enumi})}
\setcounter{enumi}{6}
\tightlist
\item
  Consider the system of scalar differential equations
\end{enumerate}

\[
\frac {d x _ {i}}{d t} = \sum_ {j, k = 1} ^ {n} c _ {i j k} x _ {j} x _ {k} \quad (1 \leqslant i \leqslant n)
\]

where the \(c_{ijk}\) are constants such that \(c_{kji} = -c_{ijk}\) . Show that the integrals of this system are defined for all values of t (note that \(\sum_{i=1}^{n} x_{i}^{2}\) is constant for every integral \(\mathbf{x} = (x_{i})\) ).

\begin{enumerate}
\def\labelenumi{\arabic{enumi})}
\setcounter{enumi}{7}
\tightlist
\item
  Let \(\mathbf{f}\) be a function defined on \(\mathrm{I} \times \mathrm{H}\) satisfying the conditions of lemma 1 of IV, p.~165, and such that for a constant \(k\) with \(0 < k < 1\) , and a point \(t_0 \in \mathrm{I}\) , one has, for all \(t \in \mathrm{I}\) and \(\mathbf{x}_1\) , \(\mathbf{x}_2 \in \mathrm{H}\) ,
\end{enumerate}

\[
\left\| \mathbf {f} (t, \mathbf {x} _ {1}) - \mathbf {f} (t, \mathbf {x} _ {2}) \right\| \leqslant \frac {k}{\left| t - t _ {0} \right|} \left\| \mathbf {x} _ {1} - \mathbf {x} _ {2} \right\|.
\]

Show that if \(\mathbf{u}\) and \(\mathbf{v}\) are two approximate solutions of the equation \(\mathbf{x}' = \mathbf{f}(t, \mathbf{x})\) , to within \(\varepsilon_1\) and \(\varepsilon_2\) respectively, defined on I, with values in H, and having the same value at the point \(t_0\) , then, for every \(t \in I\) ,

\[
\left\| \mathbf {u} (t) - \mathbf {v} (t) \right\| \leqslant \frac {\varepsilon_ {1} - \varepsilon_ {2}}{1 - k} | t - t _ {0} |.
\]

Deduce from this that th. 1 (IV, p.~171) and 2 (IV, p.~172) remain valid for the equation \(\mathbf{x}' = \mathbf{f}(t, \mathbf{x})\) under the indicated conditions.

\begin{enumerate}
\def\labelenumi{* \arabic{enumi})}
\setcounter{enumi}{8}
\tightlist
\item
  Let I be an interval in \(\mathbf{R}\) , and \(t_0\) a point of I, let S be a ball with radius \(r\) and centre \(\mathbf{x}_0\) in E, let G be the normed space of bounded maps from \(\mathrm{I} \times \mathrm{S}\) into E, the norm of such a map \(\mathbf{f}\) being the supremum \(\| \mathbf{f} \|\) of \(\| \mathbf{f}(t, \mathbf{x}) \|\) over \(\mathrm{I} \times \mathrm{S}\) . For every \(M > 0\) let \(\mathrm{G}_{\mathrm{M}}\) be the ball \(\| \mathbf{f} \| \leqslant \mathrm{M}\) in G. Let L be the subset of G formed by the Lipschitz maps of \(\mathrm{I} \times \mathrm{S}\) into E; for every function \(\mathbf{f} \in \mathrm{L} \cap \mathrm{G}_{\mathrm{M}}\) let \(\mathbf{u} = \mathrm{U}(\mathbf{f})\) be the integral of \(\mathbf{x}' = \mathbf{f}(t, \mathbf{x})\) such that \(\mathbf{u}(t_0) = \mathbf{x}_0\) , with values in S, and defined on the intersection \(\mathrm{J}_{\mathrm{M}}\) of I and the interval
\end{enumerate}

\[
\left] t _ {0} - \frac {r}{\mathbf {M}}, t _ {0} + \frac {r}{\mathbf {M}} \right[ \text {(th. 1).}
\]

\begin{enumerate}
\def\labelenumi{\alph{enumi})}
\tightlist
\item
  Let \((\mathbf{f}_{n})\) be a sequence of functions belonging to \(L \cap G_{M}\) ; if \(f_{n}\) converges uniformly on \(I \times S\) to a function f, then every cluster point (for the topology of compact convergence) of the sequence of \(\mathbf{u}_{n} = \mathrm{U}(\mathbf{f}_{n})\) in the space F of bounded maps from \(J_{M}\) into E, is an integral of \(\mathbf{x}' = \mathbf{f}(t, \mathbf{x})\) taking the value \(x_{0}\) at the point \(t_{0}\) . Conversely, every integral v of \(\mathbf{x}' = \mathbf{f}(t, \mathbf{x})\) defined on \(J_{M}\) and such that \(\mathbf{v}(t_{0}) = \mathbf{x}_{0}\) , is also an integral of an equation \(\mathbf{x}' = \mathbf{g}(t, \mathbf{x})\) , where g is Lipschitz and arbitrarily close to f in G (consider the equation
\end{enumerate}

\[
\mathbf {x} ^ {\prime} = \mathbf {f} _ {n} (t, \mathbf {x}) + \mathbf {v} ^ {\prime} (t) - \mathbf {f} _ {n} (t, \mathbf {v} (t))).
\]

\begin{enumerate}
\def\labelenumi{\alph{enumi})}
\setcounter{enumi}{1}
\tightlist
\item
  Suppose further that E is of finite dimension. Show that if \(f \in G_{M}\) satisfies the hypotheses of lemma 1, then for every \(t \in J_{M}\) the set \(H(t)\) of values at the point t of the integrals of \(\mathbf{x}' = \mathbf{f}(t, \mathbf{x})\) which take the value \(x_{0}\) at the point \(t_{0}\) is a compact and connected set (to see that \(H(t)\) is closed, use Ascoli's theorem; to see that it is connected, use a): if \(x_{1}, x_{2}\) are two points of \(H(t)\) and \(\varepsilon > 0\) is arbitrary, show that there exists a connected set \(\Phi\) of functions g belonging to \(L \cap G_{M}\) such that \(\|f - g\| \leqslant \varepsilon\) for every function \(g \in \Phi\) , and that the set of values of the functions \(U(\mathbf{g})\) at the point t is connected and contains \(x_{1}\) and \(x_{2}\) . Conclude by passing to the limit along an ultrafilter finer than the neighbourhood filter of 0 in \(R_{+}\) ).
\end{enumerate}

\begin{enumerate}
\def\labelenumi{\arabic{enumi})}
\setcounter{enumi}{9}
\tightlist
\item
  Let \(f\) be a continuous bounded real function defined on the box \(\mathrm{P}: |t - t_0| < a\) , \(|x - x_0| < b\) of \(\mathbf{R}^2\) . Let \(\mathbf{M}\) be the supremum of \(|f(t,x)|\) over \(\mathrm{P}\) , and let \(\mathrm{I} = ]t_0 - \alpha, t_0 + \alpha[\) , where \(\alpha = \inf(a, b/\mathrm{M})\) . Show that the upper and lower envelopes of the set \(\Phi\) of integrals of \(x' = f(t,x)\) defined on \(\mathrm{I}\) and taking the value \(x_0\) at the point \(t_0\) , are again integrals of \(x' = f(t,x)\) , which one calls the maximal and minimal integral of this equation corresponding to the point \((t_0, x_0)\) (remark that the set \(\Phi\) is equicontinuous and closed for the topology of uniform convergence on \(\mathrm{I}\) ).
\end{enumerate}

For every \(\tau \in I\) let \(\xi\) be the value of the minimal integral (corresponding to \((t_0, x_0)\) ) at the point \(\tau\) . Show that the minimal integral corresponding to the point \((\tau, \xi)\) is identical to the minimal integral corresponding to the point \((t_0, x_0)\) on an interval of the form \([\tau, \tau + h[\) if \(\tau > t_0\) , and of the form \(]\tau - h, \tau]\) if \(\tau < t_0\) .

Deduce from this that there exists a largest open interval \(]t_{1}, t_{2}[\) containing \(t_{0}\) and contained in \(]t_{0}-a, t_{0}+a[\) , such that the minimal integral u corresponding to \((t_{0}, x_{0})\) can be extended by continuity to \(]t_{1}, t_{2}[\) , so that at every point t of \(]t_{1}, t_{2}[\) , the value \(u(t)\) belongs to \(]x_{0}-b, x_{0}+b[\) , and that u is identical with the minimal integral corresponding to the point \((t, u(t))\) on an interval of the form \([t, t+h[\) if \(t > t_{0}\) and of the form \(]t-h, t]\) if \(t < t_{0}\) ; show further that either \(t_{1} = t_{0} - a\) (resp. \(t_{2} = t_{0} + a\) ) or that \(\lim_{t \to t_{1}} u(t) = x_{0} \pm b\) (resp. \(\lim_{t \to t_{2}} u(t) = x_{0} \pm b\) ).

\begin{enumerate}
\def\labelenumi{\arabic{enumi})}
\setcounter{enumi}{10}
\tightlist
\item
  \begin{enumerate}
  \def\labelenumii{\alph{enumii})}
  \tightlist
  \item
    On the box P: \(|t - t_{0}| < a\) , \(|x - x_{0}| < b\) let g and h be two continuous real functions such that \(g(t, x) < h(t, x)\) on P. Let u (resp. v) be an integral of \(x' = g(t, x)\) (resp. \(x' = h(t, x)\) ) such that \(u(t_{0}) = x_{0}\) (resp. \(v(t_{0}) = x_{0}\) ) defined on an interval \([t_{0}, t_{0} + c[;\) show that for \(t_{0} < t < t_{0} + c\) one has \(u(t) < v(t)\) (consider the supremum \(\tau\) of the t for which this inequality holds).
  \end{enumerate}
\end{enumerate}

\begin{enumerate}
\def\labelenumi{\alph{enumi})}
\setcounter{enumi}{1}
\item
  Let u be the maximal integral of \(x' = g(t, x)\) corresponding to the point \((t_0, x_0)\) (exerc. 10), defined on an interval \([t_0, t_0 + c[\) , with values in \(|x - x_0| < b\) . Show that in every compact interval \([t_0, t_0 + d]\) contained in \([t_0, t_0 + c[\) the minimal and maximal integrals of the equation \(x' = g(t, x) + \varepsilon\) corresponding to the point \((t_0, x_0)\) are defined once \(\varepsilon > 0\) is small enough, and converge uniformly to u when \(\varepsilon\) tends to 0 through values \textgreater{} 0.
\item
  Let g and h be two real continuous functions defined on P and such that \(g(t, x) \leqslant h(t, x)\) on P. Let \([t_0, t_0 + c[\) be an interval on which are defined an integral u of \(x' = g(t, x)\) such that \(u(t_0) = x_0\) , and the maximal integral v of \(x' = h(t, x)\) corresponding to the point \((t_0, x_0)\) . Show that \(u(t) \leqslant v(t)\) on this interval (reduce to case a) using b)).
\end{enumerate}

\begin{enumerate}
\def\labelenumi{\arabic{enumi})}
\setcounter{enumi}{11}
\tightlist
\item
  Let \(u\) be the integral of the equation \(x' = \lambda + \frac{x^2}{1 + t^2}\) equal to 0 for \(t = 0\) , and let \(J\) be the largest interval with left-hand endpoint 0 on which \(u\) is continuous.
\end{enumerate}

\begin{enumerate}
\def\labelenumi{\alph{enumi})}
\item
  Show that if \(\lambda \leqslant \frac{1}{4}\) one has \(J = [0, +\infty[\) (use exerc. 4 of IV, p.~199).
\item
  Show that if \(\lambda >\frac{1}{4}\) one has \(J = [0,a[\) where
\end{enumerate}

\[
\sinh \frac {\pi}{2 \sqrt {\lambda}} <   a <   \sinh \frac {\pi}{\sqrt {4 \lambda - 1}}
\]

(put \(x = y\sqrt{1 + t^{2}}\) and use exerc. 11).

* 13) a) Let \(\mathrm{I} = [t_0, t_0 + c]\) and let \(\omega\) be a continuous real function, \(\geqslant 0\) , defined on \(\mathrm{I} \times \mathbf{R}\) . Let \(\mathrm{S}\) be a ball with centre \(\mathbf{x}_0\) in a complete normed space \(\mathrm{E}\) , and let \(\mathbf{f}\) be a continuous map from \(\mathrm{I} \times \mathrm{S}\) into \(\mathrm{E}\) such that for any \(t \in \mathrm{I}\) , \(\mathbf{x}_1 \in \mathrm{S}\) and \(\mathbf{x}_2 \in \mathrm{S}\) one has \(\| \mathbf{f}(t, \mathbf{x}_1) - \mathbf{f}(t, \mathbf{x}_2) \| \leqslant \omega(t, \| \mathbf{x}_1 - \mathbf{x}_2 \|)\) . Let \(\mathbf{u}\) and \(\mathbf{v}\) be two integrals of \(\mathbf{x}' = \mathbf{f}(t, \mathbf{x})\) defined on \(\mathrm{I}\) , with values in \(\mathrm{S}\) , such that \(\mathbf{u}(t_0) = \mathbf{x}_1\) , \(\mathbf{v}(t_0) = \mathbf{x}_2\) ; let \(w\) be the maximal integral (exerc. 10) of \(z' = \omega(t, z)\) corresponding to the point ( \(t_0, \| \mathbf{x}_1 - \mathbf{x}_2 \|\) ), assumed to be defined on \(\mathrm{I}\) ; show that, on \(\mathrm{I}\) , one has \(\| \mathbf{u}(t) - \mathbf{v}(t) \| \leqslant w(t)\) . (Let \(w(t, \varepsilon)\) be the maximal integral of \(z' = \omega(t, z) + \varepsilon\) corresponding to the point ( \(t_0, \| \mathbf{x}_1 - \mathbf{x}_2 \|\) ), where \(\varepsilon > 0\) is sufficiently small. Show that \(\| \mathbf{u}(t) - \mathbf{v}(t) \| \leqslant w(t, \varepsilon)\) for every \(\varepsilon > 0\) , arguing by contradiction.)

\begin{enumerate}
\def\labelenumi{\alph{enumi})}
\setcounter{enumi}{1}
\tightlist
\item
  In the statement of the hypotheses of a), replace I by the interval \(I' = ]t_{0} - c, t_{0}]\) . Show that if w is the minimal integral on this interval of \(z' = \omega(t, z)\) corresponding to the point \((t_{0}, \| \mathbf{x}_{1} - \mathbf{x}_{2}\|)\) , then \(\| \mathbf{u}(t) - \mathbf{v}(t) \| \geqslant w(t)\) on \(I'\) (same method).
\end{enumerate}

* 14) a) Let \(\omega(t,z)\) be a continuous real function, \(\geqslant 0\) , defined for \(0 < t \leqslant a\) and \(z \geqslant 0\) . Suppose that \(z = 0\) is the only integral of \(z' = \omega(t,z)\) defined for \(0 < t \leqslant a\) and such that \(\lim_{t \to 0} z(t) = 0\) and \(\lim_{t \to 0} z(t)/t = 0\) . Let \(\mathrm{I} = [\mathrm{t}_0, t_0 + a[\) , let \(\mathrm{S}\) be a ball with centre \(\mathbf{x}_0\) in \(\mathrm{E}\) , and \(\mathbf{f}\) a map from \(\mathrm{I} \times \mathrm{S}\) into \(\mathrm{E}\) , such that for any \(t \in \mathrm{I}\) , \(\mathbf{x}_1 \in \mathrm{S}\) and \(\mathbf{x}_2 \in \mathrm{S}\) one has \(\|\mathbf{f}(t, \mathbf{x}_1) - \mathbf{f}(t, \mathbf{x}_2)\| \leqslant \omega(|t - t_0|, \| \mathbf{x}_1 - \mathbf{x}_2 \|)\) . Show that, on an interval with left-hand endpoint \(t_0\) contained in \(\mathrm{I}\) , the equation \(\mathbf{x}' = \mathbf{f}(t, \mathbf{x})\) can have only one solution \(\mathbf{u}\) such that \(\mathbf{u}(t_0) = \mathbf{x}_0\) . (Argue by contradiction: if \(\mathbf{v}\) is a second integral of \(\mathbf{x}' = \mathbf{f}(t, \mathbf{x})\) , bound \(\|\mathbf{u}(t) - \mathbf{v}(t)\|\) below on an interval with left-hand endpoint \(t_0\) , with the help of exerc. 13 b), and thus obtain a contradiction.)

Apply to the case where \(\omega(t, z) = k(z/t)\) with \(0 \leqslant k < 1\) (cf.~IV, p.~200, exerc. 8).

\begin{enumerate}
\def\labelenumi{\alph{enumi})}
\setcounter{enumi}{1}
\item
  The result of a) applies for \(\omega(t,z)=z/t\) ; but show by an example in this case that if u, v are two integrals to within \(\varepsilon\) of \(\mathbf{x}^{\prime}=\mathbf{f}(t,\mathbf{x})\) , equal to \(x_{0}\) at the point \(t_{0}\) , it is not possible to bound \(\|\mathbf{u}(t)-\mathbf{v}(t)\|\) above by a number depending only on t (and not on the function f). (Take for f a continuous map of \(R_{+}\times R\) into R, equal to x/t for \(t\geqslant\alpha\) and for \(0<t<\alpha\) and \(|x|\leqslant t^{2}/(\alpha-t)\) , and independent of x for the other points \((t,x)\) .)
\item
  Let \(\theta\) be a continuous real function, \(\geqslant 0\) on the interval \([0, a]\) . Show that if the integral \(\int_0^a \frac{\theta(t)}{t} dt\) converges then the result of \(a)\) applies for \(\omega(t, z) = \frac{1 + \theta(t)}{t} z\) ; on the other hand, if \(\int_0^a \frac{\theta(t)}{t} dt\) is infinite, give an example of a continuous function \(\mathbf{f}\) such that \(\| \mathbf{f}(t, \mathbf{x}_1) - \mathbf{f}(t, \mathbf{x}_2) \| \leqslant \frac{1 + \theta(|t - t_0|)}{|t - t_0|} \| \mathbf{x}_1 - \mathbf{x}_2 \|\) , but such that the equation \(\mathbf{x}' = \mathbf{f}(t, \mathbf{x})\) has infinitely many integrals equal to \(\mathbf{x}_0\) at the point \(t_0\) (method similar to that of \(b\) )).
\end{enumerate}

\begin{enumerate}
\def\labelenumi{\arabic{enumi})}
\setcounter{enumi}{14}
\item
  Let \(f\) be a real function, defined and continuous for \(|t| \leqslant a\) , \(|x| \leqslant b\) , such that \(f(t,x) < 0\) for \(tx > 0\) and \(f(t,x) > 0\) for \(tx < 0\) ; show that \(x = 0\) is the only integral of the equation \(x' = f(t,x)\) which takes the value 0 at the point \(t = 0\) (argue by contradiction).
\item
  Let E be a vector space of finite dimension, and f a bounded function on I×H satisfying the conditions of IV, p.~165, lemma 1, such that the equation \(\mathbf{x}^{\prime}=\mathbf{f}(t,\mathbf{x})\) has a unique solution u defined on I, with values in H, equal to \(x_{0}\) at the point \(t_{0}\) . Suppose further that for sufficiently large n there is an approximate integral \(u_{n}\) of \(\mathbf{x}^{\prime}=\mathbf{f}(t,\mathbf{x})\) to within 1/n, defined on I, with values in H, equal to \(x_{0}\) at the point \(t_{0}\) . Show that the sequence \((\mathbf{u}_{n})\) converges uniformly to u on every compact interval contained in I (use the fact that the sequence \((\mathbf{u}_{n})\) is equicontinuous on I).
\item
  To study the equation \(x' = \sin tx\) (IV, p.~174, Example 4) one puts u = xy and considers the solutions of the corresponding equation \(u' = \frac{u}{t} + t \sin u = \mathrm{F}(t, u)\) which are \textgreater0 for t \textgreater0 (one has \(u(0) = 0\) for every solution of this kind). Denote by \(\Gamma_k\) the curve defined by the relations \((2k - 1)\pi < u < 2k\pi\) , \(t \sin u + \frac{u}{t} = 0\) for each integer \(k \geqslant 1\) ; by \(D_k\) the open set defined by \((2k - 1)\pi < u < 2k\pi\) , \(t \sin u + \frac{u}{t} < 0\) ; and by E the complement of the union of the \(\overline{D}_k\) in the set of \((t, u)\) such that t \textgreater0 and u \textgreater0.
\end{enumerate}

\begin{enumerate}
\def\labelenumi{\alph{enumi})}
\item
  Show that every integral curve which cuts a line \(u = 2k\pi\) also cuts the line \(u = (2k + 1)\pi\) .
\item
  If an integral curve C cuts a curve \(\Gamma_{k}\) at a point \((t_{0}, u_{0})\) , then the function u is increasing for 0 \textless{} t \textless{} t \textless{} \(t_{0}\), decreasing for \(t > t_{0}\) , and, when t tends to \(+\infty\) , \(u(t)\) tends to \((2k - 1)\pi\) , the curve C remaining in \(D_{k}\) for \(t > t_{0}\) .
\item
  Show that there is no integral curve C contained in E and such that \(u(t)\) tends to \(+\infty\) when t tends to \(+\infty\) . (Form the differential equation dx/du = G(u, x) between x and u along C. If C lies entirely in E, then \(x^{2} > u\) for \(u = (2k - \frac{1}{2})\pi\) ; on the other hand, estimate \(x^{2}(u)\) by using the preceding differential equation, and obtain a contradiction.)
\end{enumerate}

\(d\)) Show that for every integer \(k\) there exists one and only one integral curve C contained in E and such that \(u(t)\) tends to \(2k\pi\) when \(t\) tends to \(+\infty\) . (On putting \(v = 2k\pi - u\) one obtains the equation \(v' = \frac{v - 2k\pi}{t} + t \sin v\) ; compare this equation to \(v' = \frac{v - 2k\pi}{t} + tv\) with \(t\) tending to \(+\infty\) and \(v\) to 0.)

\begin{enumerate}
\def\labelenumi{\arabic{enumi})}
\setcounter{enumi}{17}
\tightlist
\item
  Let E be the space of sequences \(\mathbf{x}=(x_{n})_{n\in\mathbb{N}}\) of real numbers such that \(\lim_{n\to\infty}x_{n}=0\) , endowed with the norm \(\|x\|=\sup_{n}|x_{n}|\) , which is in fact a complete normed space. For each \(\mathbf{x}=(x_{n})\in\mathrm{E}\) let \(\mathbf{y}=\mathbf{f}(\mathbf{x})\) be the element \((y_{n})\) of E defined by \(y_{n}=|x_{n}|^{\frac{1}{2}}+\frac{1}{n+1}\) ; the function f is continuous on E. Show that there is no solution of the differential equation \(\mathbf{x}^{\prime}=\mathbf{f}(\mathbf{x})\) defined on a neighbourhood of 0 in R, with values in E, and equal to 0 at the point t=0 (cf.~IV, p.~202, exerc. 11).
\end{enumerate}

\setcounter{exercisegroup}{1}
\refstepcounter{exercisegroup}
\section*{Exercises on \S\arabic{exercisegroup}}
\phantomsection
\addcontentsline{toc}{section}{Exercises on \protect\S\arabic{exercisegroup}}

\begin{enumerate}
\def\labelenumi{\arabic{enumi})}
\item
  Let E be a complete normed space over R, F a topological space, J an interval of R not reducing to a single point; let \((t, \xi) \mapsto A(t, \xi)\) be a map from \(J \times F\) into \(\mathcal{L}(E)\) , such that when \(\xi\) tends to \(\xi_{0}\) , \(A(t, \xi)\) tends uniformly to \(A(t, \xi_{0})\) in J. If \(C(t, t_{0}, \xi)\) is the resolvent of the linear equation \(d\mathbf{x}/dt = A(t, \xi)\) . x show that, for every compact interval \(K \subset J\) , \(C(t, t_{0}, \xi)\) tends uniformly to \(C(t, t_{0}, \xi_{0})\) on \(K \times K\) as \(\xi\) tends to \(\xi_{0}\) .
\item
  Let \(t \mapsto A(t)\) be a regulated map from \(J\) into \(\mathcal{L}(E)\) such that, for any two points \(s, t\) of \(J\) , \(A(s)\) and \(A(t)\) commute. Put \(B(t) = \int_{t_0}^{t} A(s) ds\) . Show that the resolvent \(C(t, t_0)\) of the equation \(d\mathbf{x}/dt = A(t)\mathbf{x}\) is equal to \(\exp(B(t))\) .
\end{enumerate}

If \(t \mapsto A_1(t)\) , \(t \mapsto A_2(t)\) are two regulated functions from J into \(\mathcal{L}(\mathrm{E})\) such that, for any two points \(s\) , \(t\) of J, \(A_1(s)\) , \(A_1(t)\) , \(A_2(s)\) , \(A_2(t)\) commute pairwise, show that

\[
\exp \left(\int_ {t _ {0}} ^ {t} \left(A _ {1} (s) + A _ {2} (s)\right) d s\right) = \exp \left(\int_ {t _ {0}} ^ {t} A _ {1} (s) d s\right) \exp \left(\int_ {t _ {0}} ^ {t} A _ {2} (s) d s\right).
\]

\begin{enumerate}
\def\labelenumi{\arabic{enumi})}
\setcounter{enumi}{2}
\tightlist
\item
  Given the two matrices
\end{enumerate}

\[
A = \left( \begin{array}{c c} 0 & 1 \\ 0 & 0 \end{array} \right), \qquad \qquad B = \left( \begin{array}{c c} 0 & 0 \\ 1 & 0 \end{array} \right)
\]

show that \(\exp(A+B)\neq\exp(A)\exp(B)\) .

\begin{enumerate}
\def\labelenumi{\arabic{enumi})}
\setcounter{enumi}{3}
\item
  If \(P\) is any automorphism in \(\mathcal{L}(\mathrm{E})\) show that \(\exp(PAP^{-1}) = P\exp(A)P^{-1}\) for every \(A \in \mathcal{L}(\mathrm{E})\) .
\item
  Show by an example that a linear equation \(d\mathbf{x} / dt = A.\mathbf{x} + e^{\alpha t}\mathbf{p}(t)\) , where \(A \in \mathcal{L}(\mathrm{E})\) is independent of \(t\) and \(\mathbf{p}\) is a polynomial (with coefficients in E) can have an integral which is equal to a polynomial of the same degree as \(\mathbf{p}\) , even when \(\alpha\) is a characteristic root of \(A\) .
\item
  Let \(f(X)\) be a polynomial of degree n with complex coefficients, and let
\end{enumerate}

\[
\frac {1}{f (\mathrm{X})} = \sum_ {j = 1} ^ {q} \sum_ {h = 1} ^ {n _ {j}} \frac {\lambda_ {j h}}{(\mathrm{X} - r _ {j}) ^ {h}}
\]

be the decomposition of the rational fraction \(1 / f(\mathrm{X})\) into simple elements (Alg., VII. 7, n° 2). Show that, for any regulated function \(b\) , the function

\[
t \mapsto \sum_ {j = 1} ^ {q} \sum_ {h = 1} ^ {n _ {j}} \lambda_ {j h} \int_ {t _ {0}} ^ {t} \frac {(t - s) ^ {h - 1}}{(h - 1) !} e ^ {r _ {j} (t - s)} b (s) d s
\]

is an integral of the equation \(f(\mathrm{D})x = b(t)\) .

\begin{enumerate}
\def\labelenumi{\arabic{enumi})}
\setcounter{enumi}{6}
\tightlist
\item
  \begin{enumerate}
  \def\labelenumii{\alph{enumii})}
  \tightlist
  \item
    Let \(t \mapsto A(t)\) be a map from an interval \(J \subset R\) into \(\mathcal{L}(E)\) , such that, for every \(x \in E\) , the map \(t \mapsto A(t).x\) of J into E is continuous. Show that \(t \mapsto \|A(t)\|\) is bounded on every compact interval \(K \subset J\) (use Gen.~Top., IX, p.~194, th. 2). Under these conditions show that, for every point \((t_{0}, \mathbf{x}_{0}) \in \mathbf{J} \times \mathbf{E}\) the equation \(d\mathbf{x}/dt = A(t).\mathbf{x} + \mathbf{b}(t)\) has one and only one solution defined on J and equal to \(x_{0}\) at the point \(t_{0}\) . Further, if one denotes this solution by \(\mathbf{u}(t, t_{0}, \mathbf{x}_{0})\) , the map \(\mathbf{x}_{0} \mapsto \mathbf{u}(t, t_{0}, \mathbf{x}_{0})\) is a bijective bicontinuous linear map \(C(t, t_{0})\) of E onto itself, which satisfies the relations (6) (IV, p.~179) and (7) (IV, p.~181); further, the map \((s, t) \mapsto C(s, t)\) is continuous from \(J \times J\) into \(\mathcal{L}(E)\) .
  \end{enumerate}
\end{enumerate}

\begin{enumerate}
\def\labelenumi{\alph{enumi})}
\setcounter{enumi}{1}
\tightlist
\item
  Take for E the space of sequences \(\mathbf{x}=(x_{n})_{n\in\mathbf{N}}\) of real numbers such that \(\lim_{n\to\infty}x_{n}=0\) , endowed with the norm \(\|x\|=\sup_{n}|x_{n}|\) , under which E is complete. For each \(t\in J=[0,1]\) , let \(A(t)\) be the linear map of E into itself such that
\end{enumerate}

\[
A (t). \mathbf {x} = \left(\frac {1}{1 + n t} x _ {n}\right) _ {n \in \mathbf {N}}.
\]

Show that \(A(t)\) satisfies the conditions of a) but that the map \(t \mapsto A(t)\) of J into \(\mathcal{L}(\mathrm{E})\) is not continuous at the point t = 0, and that the resolvent \(C(t, t_0)\) of the equation \(d\mathbf{x}/dt = A(t)\) . x is such that the map \(t \mapsto C(t, t_0)\) of J into \(\mathcal{L}(\mathrm{E})\) is not differentiable at the point t = 0.

* 8) Let G be a complete normed algebra with unit e over the field R.

\begin{enumerate}
\def\labelenumi{\alph{enumi})}
\item
  Let \(t \mapsto \mathbf{a}(t)\) be a regulated function on an interval \(J \subset R\) with values in G. Show that the integral u of the linear equation \(d\mathbf{x}/dt = \mathbf{a}(t)\mathbf{x}\) which takes the value e at a point \(t_0 \in J\) is invertible on J, and that its inverse is the solution of the equation \(d\mathbf{x}/dt = -\mathbf{x}\mathbf{a}(t)\) (if v is the integral of this last equation which takes the value e at the point \(t_0\) , consider the linear equations satisfied by uv and vu). Deduce that, for every \(x_0 \in G\) , the integral of \(d\mathbf{x}/dt = \mathbf{a}(t)\mathbf{x}\) which takes the value \(x_0\) at the point \(t_0\) is equal to \(\mathbf{u}(t)\mathbf{x}_0\) .
\item
  Let \(\mathbf{a}(t)\) , \(\mathbf{b}(t)\) be two regulated functions on J, and u and v the integrals of the equations \(d\mathbf{x}/dt = \mathbf{a}(t)\mathbf{x}\) , \(d\mathbf{x}/dt = \mathbf{x}\mathbf{b}(t)\) taking the value e at the point \(t_{0}\) . Show that the integral of the equation \(d\mathbf{x}/dt = \mathbf{a}(t)\mathbf{x} + \mathbf{x}\mathbf{b}(t)\) which takes the value \(x_{0}\) at the point \(t_{0}\) is equal to \(\mathbf{u}(t)\mathbf{x}_{0}\mathbf{v}(t)\) .
\item
  Let \(\mathbf{a}(t)\) , \(\mathbf{b}(t)\) , \(\mathbf{c}(t)\) , \(\mathbf{d}(t)\) be four regulated functions on \(\mathbf{J}\) , and \((\mathbf{u}, \mathbf{v})\) a solution of the system of two linear equations
\end{enumerate}

\[
\frac {d \mathbf {x}}{d t} = \mathbf {a} (t) \mathbf {x} + \mathbf {b} (t) \mathbf {y}, \quad \frac {d \mathbf {y}}{d t} = \mathbf {c} (t) \mathbf {x} + \mathbf {d} (t) \mathbf {y}.
\]

Show that, if \(\mathbf{v}\) is invertible on \(J\) , then \(\mathbf{w} = \mathbf{u}\mathbf{v}^{-1}\) is an integral of the equation \(d\mathbf{z}/dt = \mathbf{b}(t) + \mathbf{a}(t)\mathbf{z} - \mathbf{z}\mathbf{d}(t) - \mathbf{z}\mathbf{c}(t)\mathbf{z}\) (the ``Ricatti equation''); and conversely. Deduce that every integral of this last equation on \(J\) , taking the value \(\mathbf{x}_0\) at the point \(t_0\) , is of the form \((A(t)\mathbf{x}_0 + B(t)\mathbf{e})(C(t)\mathbf{x}_0 + D(t)\mathbf{e})^{-1}\) , where \(A(t), B(t), C(t)\) and \(D(t)\) are continuous maps from \(J\) into \(\mathcal{L}(E)\) satisfying the identity \(U.(xy) = (U.x)y\) .

\begin{enumerate}
\def\labelenumi{\alph{enumi})}
\setcounter{enumi}{3}
\tightlist
\item
  Let \(w_{1}\) be an integral of \(d\mathbf{z}/dt = \mathbf{b}(t) + \mathbf{a}(t)\mathbf{z} - \mathbf{z}\mathbf{d}(t) - \mathbf{z}\mathbf{c}(t)\mathbf{z}\) on J; show that if w is another integral of this equation such that \(w - w_{1}\) is invertible on J then \(w - w_{1}\) can be expressed in terms of integrals of the equations \(d\mathbf{x}/dt = -(\mathbf{d} + \mathbf{c}\mathbf{w}_{1})\mathbf{x}\) and \(d\mathbf{x}/dt = \mathbf{x}(\mathbf{a} - \mathbf{w}_{1}\mathbf{c})\) (consider the equation satisfied by \((\mathbf{w} - \mathbf{w}_{1})^{-1}\) , and use b)).
\end{enumerate}

\begin{enumerate}
\def\labelenumi{\arabic{enumi})}
\setcounter{enumi}{8}
\tightlist
\item
  Let \(y_{k}\) ( \(1 \leqslant k \leqslant n\) ) be \(n\) real functions defined on an interval \(I \subset \mathbf{R}\) and having a continuous \((n - 1)^{th}\) derivative.
\end{enumerate}

\begin{enumerate}
\def\labelenumi{\alph{enumi})}
\item
  Show that if the \(n\) functions \(y_{k}\) are linearly dependent then the matrix \((y_{k}^{(h)}(t))_{0\leqslant h\leqslant n - 1, 1\leqslant k\leqslant n}\) has rank \(< n\) at every point \(t\in I\) .
\item
  Conversely, if for every \(t \in I\) the matrix \((y_{k}^{(h)}(t))_{0 \leqslant h \leqslant n-1, 1 \leqslant k \leqslant n}\) has rank \textless{} n, show that in every nonempty open interval \(J \subset I\) there exists a nonempty open interval \(U \subset J\) such that the restrictions of the \(y_{k}\) to U are linearly dependent (if p is the smallest of the numbers q \textless{} n such that the Wronskians of any q of the functions \(y_{k}\) are identically zero on J, consider a point \(a \in J\) where the Wronskian of p - 1 of the functions \(y_{k}\) is not zero, and show that p of the functions \(y_{k}\) are integrals of a linear equation of order p - 1 on a neighbourhood of a).
\item
  Put \(y_{1}(t) = t^{2}\) for \(t \in \mathbf{R}\) , \(y_{2}(t) = t^{2}\) for \(t \geqslant 0\) , \(y_{2}(t) = -t^{2}\) for \(t < 0\) ; show that \(y_{1}\) and \(y_{2}\) have continuous derivatives on \(\mathbf{R}\) and that \(y_{1}y_{2}' - y_{2}y_{1}' = 0\) , but that \(y_{1}\) and \(y_{2}\) are not linearly dependent on \(\mathbf{R}\) .
\end{enumerate}

* 10) Let \(t \mapsto X(t)\) be a map from an interval \(\mathbf{J} \subset \mathbf{R}\) into the space of complex matrices of order \(n\) . Suppose that the derivative \(t \mapsto X'(t)\) exists and is continuous on \(\mathbf{J}\) , and is such that \(X(t)X'(t) = X'(t)X(t)\) for every \(t \in \mathbf{J}\) .

\begin{enumerate}
\def\labelenumi{\alph{enumi})}
\item
  Suppose further that for each \(t \in J\) the eigenvalues of \(X(t)\) are distinct. Show that then there exists a constant invertible matrix \(P_{0}\) such that \(P_{0}X(t)P_{0}^{-1} = D(t)\) , where \(D(t)\) is a diagonal matrix, so that \(X(t_{1})\) and \(X(t_{2})\) commute for every pair of values \(t_{1}, t_{2}\) of t in J. (Write \(X(t) = P(t)D(t)P(t)^{-1}\) on a neighbourhood of each point of J and form the differential equation satisfied by \(P(t)\) .)
\item
  The matrix
\end{enumerate}

\[
X (t) = \left( \begin{array}{c c c} t ^ {2} & t ^ {3} & t ^ {4} \\ - 2 t & - 2 t ^ {2} & - 2 t ^ {3} \\ 1 & t & t ^ {2} \end{array} \right)
\]

commutes with its derivative, but the conclusion of a) does not hold.*

\section*{Historical Note}
\phantomsection
\addcontentsline{toc}{section}{Historical Note}

(N.B. The roman numerals refer to the Bibliography at the end of this note.)

As we have seen (Historical Note to Chapters I-II-III), the problems that led to the integration of differential equations were among the first to be treated by the founders of the infinitesimal calculus of the XVII \(^{th}\) century (notably Descartes and Barrow). Since then the theory of differential equations has not ceased to tax the sagacity of mathematicians, and to be a favoured terrain for the application of the most varied methods of Analysis; the questions that it provokes are very far from being all solved, and the interest that attaches is accordingly so sustained that it constitutes one of the most permanent and fruitful points of contact between mathematics and the experimental sciences: these latter often find a valuable help there, and in exchange constantly provide new problems.

Of the numerous chapters which a modern complete study of differential equations should encompass, we have not wished to expound here but the two most elementary, treating existence theorems and linear equations, the variable being assumed to be real. This is why we also limit our brief historical account to these two aspects. From the beginning of the XVIII \(^{th}\) century mathematicians were convinced that the ``general'' integral of a differential equation of order n depends on n arbitrary constants, and that ``in general'' there exists one and only one integral with given values for the function and its first n - 1 derivatives at a given value \(x_{0}\) of the variable: a conviction which they justified by the process (going back to Newton) of calculating one by one the coefficients of the Taylor expansion of the solution at the point \(x_{0}\) , with the help of the differential equation itself and the first n coefficients. But up until Cauchy no one had studied the convergence of the series so obtained, nor proved that its sum was a solution of the differential equation; and of course, there was no question of nonanalytic differential equations. Among the various methods invented by Cauchy to prove the existence of integrals of differential equations, that which we have followed ((IV) and (IV bis)), generalized a little later by Lipschitz, is particularly interesting for the case of nonanalytic equations and approximation of integrals.

The linear differential equations were among the first to draw attention. Leibniz and Jakob Bernoulli integrated the first order linear equation by two quadratures ((I), v. II, p.~731); the general integral of the arbitrary linear equation with constant coefficients and arbitrary right hand side was obtained by Euler ((II), p.~296-354); d'Alembert solved the linear systems with constant coefficients in the same way. A little later, Lagrange (III) treated the general theory of linear equations of order \(n\) , recognised that the integral of the homogeneous equation is a linear function of the \(n\) constants of integration, introduced the adjoint equation, discovered the reduction of the order of the homogeneous equation when one knows particular solutions, and the method of variation of parameters to solve the inhomogeneous equation. The obscure points in Lagrange's theory (notably in what concerns the linear independence of the integrals) were clarified by Cauchy, whose exposition (IV bis) has remained quasi definitive, apart from the improvements in detail provided by matrix notation and the theory of elementary divisors.

\section*{Bibliography}
\phantomsection
\addcontentsline{toc}{section}{Bibliography}

\begin{enumerate}
\def\labelenumi{(\Roman{enumi})}
\item
  JAKOB BERNOULLI, Opera, 2 vol., Genève (Cramer-Philibert), 1744.
\item
  L. EULER, Opera Omnia: Institutiones calculi integralis, (1) t. XII, Leipzig-Berlin (Teubner), 1914.
\item
  J.-L. LAGRANGE, Œuvres, Paris (Gauthier-Villars), 1867-1890: a) Solution de différents problèmes de calcul intégral, t. I, p.~471; b) Sur le mouvement des nœuds des orbites planétaires, t. IV, p.~111.
\item
  A.-L. CAUCHY, Œuvres complètes, (2), t. XI, Paris (Gauthier-Villars), 1913, p.~399 (= Exercices d'Analyse, Paris, 1840, t. I, p.~327).
\end{enumerate}

(IV bis) A.-L. CAUCHY, dans Leçons de calcul différentiel et de calcul intégral, rédigées principalement d'après les méthodes de M. A.-L. Cauchy, par l'abbé Moigno, t. II, Paris, 1844.

\chapter{Local study of functions}

\section{COMPARISON OF FUNCTIONS ON A FILTERED SET}

Let E be a set filtered by a filter with base F (Gen.~Top., I, p.~58); in this chapter we shall consider functions whose domain of definition is a subset of E belonging to the filter base F (the subset depending on the function) and which take their values either in the field of real numbers R, or, more generally, in a normed vector space over a valued field (Gen.~Top., IX, p.~170).

In applications, E will most often be a subspace of the real space \(R^{n}\) , or the extended line \(\overline{R}\) , and \(F\) will be the trace on E of the filter of neighbourhoods of a closure point of E, or else the filter of complements of relatively compact subsets of E (``neighbourhoods of the point at infinity'').

It will not in general suffice to know that such a function tends to a given limit along \(\mathfrak{F}\) to be able to treat all the problems of ``passage to the limit along \(\mathfrak{F}\)'' for expressions formed from this function.

For example, when the real variable \(x\) tends to \(+\infty\) the three functions \(x\) , \(x^2\) and \(\sqrt{x}\) all tend to \(+\infty\) , but, of the expressions

\[
(x + 1) ^ {2} - x ^ {2}, \qquad (x + 1) - x, \quad \sqrt {x + 1} - \sqrt {x},
\]

the first tends to \(+\infty\) , the second to 1, the third to 0.

It is thus important to know not just the limit value of a function along \(\mathfrak{F}\) (when this limit exists) but also the ``manner'' in which the function approaches its limit; in other words, one is led to classify the set of functions which tend to the same limit.
